%% ============================================================================
%% main_restructured.tex --- FOUR-QUESTION RESTRUCTURE VARIANT (2026-08-13)
%% Diff against main.tex before adopting. All prose is unchanged and yours
%% except short transition passages marked "% [restructure:" and two new
%% table skeletons. Moves: the Spearman-mass and Spearman-profile analyses
%% and the 1P field-response subsection now form Sec. 5 ("Fidelity across
%% the Parameter Space", with an empty Redshift Dependence stub); Sec. 4 is
%% retitled "Fidelity at Fixed Parameters"; "The BIND Sandbox" is retitled
%% "Emergent Joint Structure"; Q-tags added so Secs. 4-7 answer the four
%% introduction questions in order. If adopted, rename this file main.tex.
%% ============================================================================
%% Beginning of file 'sample701.tex'
%%
%% Version 7.0.1. Created May 2025.
%% Version 7. Created January 2025.  
%%
%% AASTeX v7+ calls the following external packages:
%% times, hyperref, ifthen, hyphens, longtable, xcolor, 
%% bookmarks, array, rotating, ulem, and lineno 
%%
%% RevTeX is no longer used in AASTeX v7+.
%%
\documentclass[linenumbers,trackchanges, twocolumn]{aastex701}
%%
%% This initial command takes arguments that can be used to modify easily 
%% the output of the compiled manuscript. Any combination of arguments can be 
%% invoked like this:
%%
%% \documentclass[argument1,argument2,argument3,...]{aastex701}
%%
%% Six of the arguments are typestting options. They are:
%%
%%  twocolumn   : two text columns, 10 point font, single spaced article.
%%                This is the most compact and represent the final published
%%                derived PDF copy of the accepted manuscript from the publisher
%%  default     : one text column, 10 point font, single spaced (default).
%%  manuscript  : one text column, 12 point font, double spaced article.
%%  preprint    : one text column, 12 point font, single spaced article.  
%%  preprint2   : two text columns, 12 point font, single spaced article.
%%  modern      : a stylish, single text column, 12 point font, article with
%% 		  wider left and right margins. This uses the Daniel
%% 		  Foreman-Mackey and David Hogg design.
%%
%% Note that you can submit to the AAS Journals in any of these 6 styles.
%%
%% There are other optional arguments one can invoke to allow other stylistic
%% actions. The available options are:
%%
%%   astrosymb    : Loads Astrosymb font and define \astrocommands. 
%%   tighten      : Makes baselineskip slightly smaller, only works with 
%%                  the twocolumn substyle.
%%   times        : uses times font instead of the default.
%%   linenumbers  : turn on linenumbering. Note this is mandatory for AAS
%%                  Journal submissions and revisions.
%%   trackchanges : Shows added text in bold.
%%   longauthor   : Do not use the more compressed footnote style (default) for 
%%                  the author/collaboration/affiliations. Instead print all
%%                  affiliation information after each name. Creates a much 
%%                  longer author list but may be desirable for short 
%%                  author papers.
%% twocolappendix : make 2 column appendix.
%%   anonymous    : Do not show the authors, affiliations, acknowledgments,
%%                  and author contributions for dual anonymous review.
%%  resetfootnote : Reset footnotes to 1 in the body of the manuscript.
%%                  Useful when there are a lot of authors and affiliations
%%		    in the front matter.
%%   longbib      : Print article titles in the references. This option
%% 		    is mandatory for PSJ manuscripts.
%%
%% Since v6, AASTeX has included \hyperref support. While we have built in 
%% specific %% defaults into the classfile you can manually override them 
%% with the \hypersetup command. For example,
%%
%% \hypersetup{linkcolor=red,citecolor=green,filecolor=cyan,urlcolor=magenta}
%%
%% will change the color of the internal links to red, the links to the
%% bibliography to green, the file links to cyan, and the external links to
%% magenta. Additional information on \hyperref options can be found here:
%% https://www.tug.org/applications/hyperref/manual.html#x1-40003
%%
%% The "bookmarks" has been changed to "true" in hyperref
%% to improve the accessibility of the compiled pdf file.
%%
%% If you want to create your own macros, you can do so
%% using \newcommand. Your macros should appear before
%% the \begin{document} command.
%%
\newcommand{\vdag}{(v)^\dagger}
\newcommand\aastex{AAS\TeX}
\newcommand\latex{La\TeX}
\usepackage{bm}
\usepackage{amsmath}
%%%%%%%%%%%%%%%%%%%%%%%%%%%%%%%%%%%%%%%%%%%%%%%%%%%%%%%%%%%%%%%%%%%%%%%%%%%%%%%%
%%
%% The following section outlines numerous optional output that
%% can be displayed in the front matter or as running meta-data.
%%
%% Running header information. A short title on odd pages and 
%% short author list on even pages. Note that this
%% information may be modified in production.
\shorttitle{BIND}
\shortauthors{Lee et al.}
%%
%% Include dates for submitted, revised, and accepted.
%%\received{February 1, 2025}
%%\revised{March 1, 2025}
%%\accepted{\today}
%%
%% Indicate AAS Journal the manuscript was submitted to.
%%\submitjournal{PSJ}
%% Note that this command adds "Submitted to " the argument.
%%
%% You can add a light gray and diagonal water-mark to the first page 
%% with this command:
%% \watermark{text}
%% where "text", e.g. DRAFT, is the text to appear.  If the text is 
%% long you can control the water-mark size with:
%% \setwatermarkfontsize{dimension}
%% where dimension is any recognized LaTeX dimension, e.g. pt, in, etc.
%%%%%%%%%%%%%%%%%%%%%%%%%%%%%%%%%%%%%%%%%%%%%%%%%%%%%%%%%%%%%%%%%%%%%%%%%%%%%%%%
%%
%% Use this command to indicate a subdirectory where figures are located.
%%\graphicspath{{./}{figures/}}
%% This is the end of the preamble.  Indicate the beginning of the
%% manuscript itself with \begin{document}.

\begin{document}

\title{BIND (Baryonic INpainting with Deep learning): A Field-level Emulator for Galaxy Groups and Clusters}

%% A significant change from AASTeX v6+ is in the author blocks. Now an email
%% address is required for each author. This means that each author requires
%% at least one of the following:
%%
%% \author
%% \affiliation
%% \email
%%
%% If these three commands are not available for each author, the latex
%% compiler will issue an error and if you force the latex compiler to continue,
%% it will generate an incomplete pdf.
%%
%% Multiple \affiliation commands are allowed and authors can also include
%% an optional \altaffiliation to indicate a status, i.e. Hubble Fellow. 
%% while affiliations are indexed as footnotes, altaffiliations are noted with
%% with a non-numeric footnote that is set away from the numeric \affiliation 
%% footnotes. NOTE that if an \altaffiliation command is used it must 
%% come BEFORE the \affiliation call, right after the \author command, in 
%% order to place the footnotes in the proper location. Because non-numeric
%% symbols are used, \altaffiliation should be used sparingly.
%%
%% In v7+ the \author command takes an optional argument which provides 
%% additional metadata about the author. Authors can provide the 16 digit 
%% ORCID, the surname (family or last) name, the given (first or fore-) name, 
%% and a name suffix, e.g. "Jr.". The syntax is:
%%
%% \author[orcid=0000-0002-9072-1121,gname=Gregory,sname=Schwarz]{Greg Schwarz}
%%
%% This name metadata in not shown, it is only for parsing by the peer review
%% system so authors can be more easily identified. This name information will
%% also be sent to the publisher so they can include it in the CROSSREF 
%% metadata. Including an orcid will hyperlink the author name to the 
%% author's ORCID page. Note that  during compilation, LaTeX will do some 
%% limited checking of the format of the ID to make sure it is valid. If 
%% the "orcid-ID.png" image file is  present or in the LaTeX pathway, the 
%% ORCID icon will appear next to the authors name.
%%
%% Even though emails are now required for each author, the \email does not
%% produce output in the compiled manuscript unless the optional "show" command
%% is used. For example,
%%
%% \email[show]{greg.schwarz@aas.org}
%%
%% All "shown" emails are show in the bottom left of the first page. Due to
%% space constraints, only a few emails should be shown. 
%%
%% To identify a corresponding author, use the \correspondingauthor command.
%% The command appends "Corresponding Author: " to the argument it appears at
%% the bottom left of the first page like the output from \email. 

\author[orcid=0000-0002-2318-3087,sname='M.E. Lee']{Max E. Lee}
\affiliation{Department of Astronomy, Columbia University, MC 5246, 538 West 120th Street, New York, NY 10027, USA}
\email[show]{max.e.lee@columbia.edu}  

\author[0000-0002-3185-1540]{Shy Genel}
\affiliation{Center for Computational Astrophysics, Flatiron Institute, 162 Fifth Ave, New York, NY, 10010, USA}
\affiliation{Columbia Astrophysics Laboratory, Columbia University, 550 West 120th Street, New York, NY 10027, USA}
\email{sgenel@flatironinstitute.org}


\author[0000-0003-3633-5403]{Zolt\'an Haiman}
\affiliation{Department of Astronomy, Columbia University, MC 5246, 538 West 120th Street, New York, NY 10027, USA}
\affiliation{Department of Physics, Columbia University, MC 5255, 538 West 120th Street, New York, NY 10027, USA}
\affiliation{Institute of Science and Technology Austria, Am Campus 1, Klosterneuburg 3400 Austria}
\email{Zoltan.Haiman@ista.ac.at}

\author[0000-0003-2630-9228]{Greg L. Bryan}
\affiliation{Department of Astronomy, Columbia University, MC 5246, 538 West 120th Street, New York, NY 10027, USA}
\affiliation{Center for Computational Astrophysics, Flatiron Institute, 162 Fifth Ave, New York, NY, 10010, USA}
\email{gbryan@columbia.edu}


%% Use the \collaboration command to identify collaborations. This command
%% takes an optional argument that is either a number or the word "all"
%% which tells the compiler how many of the authors above the command to
%% show. For example "\collaboration[all]{(DELVE Collaboration)}" wil include
%% all the authors above this command.
%%
%% Mark off the abstract in the ``abstract'' environment. 
\begin{abstract}
Baryonic feedback is the leading astrophysical systematic for upcoming
weak-lensing surveys, yet the tools used to model it are analytic,
deterministic, and calibrated almost entirely on two-point statistics. We
introduce \textsc{BIND} (Baryonic INpainting with Deep learning), a conditional
flow matching model that learns the field-level mapping from dark-matter-only
halos to their hydrodynamical counterparts. Trained on matched halo patches
from the 1{,}024 paired simulations of the CAMELS SB35 suite, \textsc{BIND}
samples the projected dark matter, gas, and stellar mass fields given the
dark-matter-only patch, the redshift, and the full 35-dimensional space of
$\Lambda$CDM and IllustrisTNG subgrid parameters, imposing no analytic form on
any component. Across three held-out evaluation
suites, \textsc{BIND} recovers integrated masses without bias --- to better
than $1\%$ for dark matter and total mass, ${\sim}1\%$ for gas, and ${\sim}5\%$
for stars --- reproduces azimuthally averaged profiles to $\lesssim10\%$ at all
radii, and matches the distributions of projected halo shapes. The learned
parameter dependence is correct at every level we test: rank correlations
between the generated fields and the subgrid parameters match the truth for
both integrated masses and radial profiles, and the field-level response to
individual parameter variations is recovered in sign and morphology. Halo mass is never supplied as conditioning, yet the
baryon fraction, the inter-component scaling relations, and the joint
covariance of their residuals all emerge correctly. Applied halo by halo to an
$N$-body volume, \textsc{BIND} reproduces the projected matter power spectrum
suppression to the accuracy ceiling set by halo replacement itself, in minutes
on a single GPU.
\end{abstract}

%% Keywords should appear after the \end{abstract} command. 
%% The AAS Journals now uses Unified Astronomy Thesaurus (UAT) concepts:
%% https://astrothesaurus.org
%% You will be asked to selected these concepts during the submission process
%% but this old "keyword" functionality is maintained in case authors want
%% to include these concepts in their preprints.
%%
%% You can use the \uat command to link your UAT concepts back its source.
\keywords{\uat{Galaxies}{573} --- \uat{Cosmology}{343}}

%% From the front matter, we move on to the body of the paper.
%% Sections are demarcated by \section and \subsection, respectively.
%% Observe the use of the LaTeX \label
%% command after the \subsection to give a symbolic KEY to the
%% subsection for cross-referencing in a \ref command.
%% You can use LaTeX's \ref and \label commands to keep track of
%% cross-references to sections, equations, tables, and figures.
%% That way, if you change the order of any elements, LaTeX will
%% automatically renumber them.

\section{Introduction} 

Exploiting the full statistical power of future surveys requires accurate modeling of the matter distribution on small, highly non-linear scales. Baryonic processes --- including feedback from active galactic nuclei (AGN) and supernovae, radiative cooling, and star formation --- have previously been a subdominant effect in cosmological analyses, but for upcoming surveys, they will become an important source of systematic uncertainty, and it is necessary to develop methods to model them \citep[e.g.][]{Jing-2006,Chisari-2019,Schneider-2019}. In response to this demand, an industry of methods for baryonic accounting in cosmological analyses has emerged.

The most straightforward and arguably cost-efficient treatment of baryonic effects is to remove small-scale data directly, a strategy adopted in several recent weak-lensing analyses, including DES, KiDS, and HSC \citep[e.g.][]{Secco-2022,Li-2023,Dalal-2023}. Such cuts discard precisely the non-linear scales that baryons dominate at but which have also been shown to contain a wealth of cosmological information, ultimately leading to a non-negligible loss in cosmology constraining power \citep[e.g.][]{VillaescusaNavarro-2021a}.
% TODO(Max): "Horowitz2017" had no identifiable match --- restore the
% intended second reference for the small-scale-information claim.

The alternative extreme is to model baryonic effects directly with cosmological hydrodynamical simulations that self-consistently evolve dark matter and baryons and, when appropriately calibrated, serve as ideal testing grounds for baryonic effects and cosmological analyses. However, these simulations, such as IllustrisTNG, EAGLE, SIMBA, BAHAMAS, or FLAMINGO \citep[e.g.][]{Vogelsberger-2014,Schaye-2015,Pillepich-2018,Dave-2019,Schaye-2023}, are computationally expensive, and the galaxy-formation physics below the resolution scale is implemented through uncertain subgrid prescriptions that differ from code to code \citep{vanDaalen-2011,Vogelsberger-2020,VillaescusaNavarro-2021}. The high computational cost severely limits their use in forward-modeling pipelines that require sampling over cosmological and galaxy-formation parameters.

To reduce the computational cost of hydrodynamical simulations while avoiding small-scale data cuts, semi-analytic baryonic correction models (BCMs) have been proposed and used. BCMs modify dark-matter-only (DMO) simulations, which are far cheaper than hydrodynamical simulations, by applying physically motivated empirical prescriptions to halo density profiles \citep{Schneider-2015,Schneider-2019,Arico-2020}.
% TODO(Max): "SchneiderBridle2010" (Schneider & Bridle 2010 is an IA halo
% model, not a BCM) --- replace with the intended early-BCM reference
% (Semboloni et al. 2011?) if one was meant.
In these models, a DMO halo profile is decomposed into components such as cold bound gas, ejected gas, and relaxed dark matter, each with an analytic form. The A20--BCM introduced by \citet{Arico-2020} has been shown to reproduce the 3D matter power spectrum suppression due to baryons at the percent level across several state-of-the-art hydrodynamical simulations. Building on this, the BACCO project constructed a neural-network-based emulator over a multi-dimensional baryonic and cosmological parameter space, enabling fast evaluation of BCM predictions for $P(k)$ \citep{Arico-2021}. Similar to BCM models are halo-model approaches such as HMcode-2020, which introduces a small number of baryonic feedback parameters into an augmented halo model to match the matter power spectrum to percent-level accuracy \citep{Mead-2021}. Like BCMs, these models provide an efficient description of power-spectrum-level baryonic effects, but they do not directly specify the underlying field-level redistribution of mass within and around halos. As a result, their impact on higher-order and non-Gaussian statistics must be calibrated a posteriori and may not faithfully represent the true spatial structure induced by baryons.

In parallel, the more recent component-wise baryonification (BFC) framework of \citet{Schneider-2025} extended previous BCMs by assigning each DMO particle from an $N$-body simulation associated dark matter, gas, and stellar components, which are displaced individually and assigned thermodynamic properties using hydrostatic and ideal gas equations. The BFC model was validated against the FLAMINGO and IllustrisTNG simulation suites, achieving $\lesssim 2\%$ agreement in the matter power spectrum suppression up to $k = 5~h\,\mathrm{Mpc}^{-1}$ across a range of feedback prescriptions. In a companion paper, \citet{Kovac-2025} applied the BFC framework to jointly fit kinematic Sunyaev-Zel'dovich (kSZ) observations from the Atacama Cosmology Telescope (ACT) and X-ray gas fractions from eROSITA, finding that current multiwavelength data prefer a stronger feedback model than assumed in most hydrodynamical simulations, with a resulting power spectrum suppression of $20$--$25\%$ at $k = 5~h\,\mathrm{Mpc}^{-1}$. This is consistent with recent findings by \citet{Hadzhiyska-2024} and \citet{Bigwood-2024, Siegel-2025}, who also find that standard feedback models appear to underpredict the level of baryon suppression inferred from observations.
% TODO(Max): identify the intended "lau" reference (E. Lau et al. 2025
% group thermodynamics?) and restore it here.

BCMs, halo-model approaches, and BFCs have been calibrated and tested almost exclusively on two-point statistics. The first steps beyond this regime, such as joint fits to the power spectrum and bispectrum, have shown that while BCMs can reach $\lesssim 3\%$ agreement for the bispectrum, discrepancies remain on small scales and for certain triangle configurations \citep{Arico-2021a}. Moreover, when tested on non-Gaussian statistics such as peak counts, BCMs can fail more dramatically: \citet{Lee-2023} found that while the A20--BCM reproduces weak-lensing peak counts at the percent level for peaks with signal-to-noise $S/N > 4$, it systematically underpredicts the number of the highest peaks, particularly for deep, wide-area surveys such as LSST and \textit{Euclid}.

More recently, several field-level and map-level approaches have been developed to go beyond analytic halo profiles. \citet{Sharma-2024} proposed a model for baryonic effects where DMO simulation modes are multiplied by an isotropic Fourier-space transfer function, $\sqrt{P_{\rm hydro}(k)/P_{\rm DMO}(k)}$. This approach is motivated by the observation that the cross-correlation coefficient between matched hydrodynamical and $N$-body simulations remains close to unity down to small scales \citep{Sharma-2024}, and has been shown to be accurate for two-point statistics across thousands of simulations. However, because the transfer function is isotropic in $k$-space, this approach cannot capture anisotropic, halo-centric mass redistribution, the locations that are sourcing the baryonic effects to begin with. With advances in machine learning, several other approaches have been explored. \citet{Chadayammuri-2023} used a U-Net to map DMO cluster fields to gas density, temperature, and X-ray observables, reproducing observed mass-observable scaling relations for massive clusters. Similar halo-based approaches have used machine learning to infer baryonic properties such as gas mass and temperature from dark-matter halo properties in simulations such as EAGLE \citep[e.g.][]{Moews-2021}. Diffusion models have also been explored for cosmological fields: Cosmo-FOLD, introduced by \citet{Mishra-2026}, uses a latent diffusion model to generate 3D dark-matter and gas-density fields conditioned on lower-resolution or partial information, achieving $\lesssim 10\%$ accuracy in power spectra up to $k \lesssim 5~h\,\mathrm{Mpc}^{-1}$. These methods represent an important step toward generative, field-level emulation of baryonic effects but operate on full volumes without explicit halo-level conditioning, making it non-trivial to connect to the actual halo-level observables we know to be associated with these baryonic effects.

Despite this progress, several important gaps remain. First, BCMs and BFCs, while halo-centric, rely on fixed analytic profile prescriptions and are not generative: they cannot sample from a posterior over baryonic field configurations at fixed dark matter and parameters. Transfer-function approaches are field-level but isotropic in $k$-space and similarly non-generative. Map-level ML approaches operate on full volumes rather than individual halos, losing the halo-centric resolution needed to capture detailed mass redistribution within and around groups and clusters. No existing method simultaneously (i) operates at the field level in a halo-centric manner, (ii) is conditioned on the full galaxy-formation parameter space of a state-of-the-art hydrodynamical model, and (iii) is stochastically generative, enabling proper posterior sampling over baryonic field realizations at fixed input. Second, higher-order and non-Gaussian weak-lensing statistics, such as peak and minimum counts, Minkowski functionals, and field-level summaries, have been shown to add substantial cosmological information beyond the power spectrum, but existing baryonic correction tools have not been systematically validated for these observables \citep{Lee-2023,Lee-2026a}. Third, the interplay between galaxy-formation model parameters and field-level matter redistribution is poorly understood: the 30-dimensional IllustrisTNG parameter space exhibits complex, nonlinear degeneracies (for example, between supernova and AGN feedback), which analytic prescriptions are potentially ill-equipped to capture \citep{Lee-2024}.
Table~\ref{tab:capabilities} summarizes where existing approaches stand on these axes.
% [restructure: pointer sentence + table skeleton are new --- CHECK rows and
% wording; characterizations follow the text of this section.]

% [restructure: new table skeleton --- CHECK rows and wording; the
% characterizations follow the text of this section.]
\begin{deluxetable*}{lccccc}
\tabletypesize{\footnotesize}
\tablecaption{Capabilities of baryonic-modeling approaches across the gaps
identified in this section.\label{tab:capabilities}}
\tablewidth{0pt}
\tablehead{\colhead{Approach} & \colhead{Halo-centric} & \colhead{Field-level} &
           \colhead{Generative} & \colhead{Param.-cond.} &
           \colhead{Multi-channel}}
\startdata
BCM / halo model$^{a}$ & yes & no (profiles) & no & BCM params & no \\
BFC$^{b}$ & yes & particle-level & no & BFC params & analytic \\
Isotropic transfer function$^{c}$ & no & yes & no & suite params & no (total matter) \\
Volume-level generative$^{d}$ & no & yes & yes & no & DM, gas \\
Halo image-to-image$^{e}$ & yes & yes & no & fixed physics & gas, X-ray \\
\textsc{BIND} (this work) & yes & yes (2D patches) & yes & 35 (cosmo + TNG) & DM, gas, stars
\enddata
\tablecomments{$^{a}$\citet{Schneider-2015, Arico-2020, Mead-2021};
$^{b}$\citet{Schneider-2025}; $^{c}$\citet{Sharma-2024};
$^{d}$\citet{Mishra-2026}; $^{e}$\citet{Chadayammuri-2023}.}
\end{deluxetable*}

In this work, we introduce \emph{BIND} (Baryonic INpainting with Deep learning), a conditional flow matching model that learns the mapping from dark-matter-only halos to their hydrodynamical counterparts at the field level. BIND is trained on matched halo patches from hydrodynamical and dark-matter-only simulations across the CAMELS IllustrisTNG suite of $50\,{\rm Mpc}/h$ simulations and is conditioned on the full set of 35 cosmological and astrophysical parameters from $\Lambda$CDM and the IllustrisTNG galaxy formation model. For each halo, BIND takes as input a local patch of the DMO density field centered on the halo and outputs simultaneous multi-channel predictions for the dark matter, gas, and stellar mass distributions, without imposing analytic functional forms on any component. In a companion paper \citep{Lee-2026c}, \textsc{BIND} is further extended to output field-level thermodynamic quantities --- temperature, pressure, and entropy maps --- the key quantities used to gauge feedback at the halo level. Because the model is generative, it naturally provides samples from the posterior distribution of baryonic fields given a DMO realization and parameter vector. 

By construction, BIND is halo-centric, field-level, and parameter-conditioned. This allows it to (i) capture anisotropic mass redistribution in and around halos, (ii) learn the dependence of different mass components across the galaxy-formation parameter space, and (iii) preserve higher-order information in the output fields, enabling direct application to downstream statistics. By applying BIND patch-by-patch to the halos of an arbitrary $N$-body simulation, one can baryonify large cosmological volumes across the IllustrisTNG parameter space in minutes on a single GPU --- orders of magnitude faster than re-running the hydrodynamical simulation --- making the method suitable for field-level analysis pipelines.

In this paper, we introduce BIND and provide a comprehensive suite of validations exploring its ability to generate accurate halo fields with correct parameter dependencies. We are primarily interested in answering:
\begin{enumerate}
    \item Can a conditional flow matching model learn the mapping from dark-matter-only to hydrodynamical fields at the halo level, conditioned on the IllustrisTNG galaxy-formation parameter space?
    \item Does BIND learn the appropriate astrophysical parameter dependencies for the generated hydrodynamical fields?
    \item Is BIND able to learn the correlations between field-level observables, and how they change as a function of astrophysical parameters?
    \item Can BIND be deployed as a practical baryonification engine for two-point statistics such as the projected matter power spectrum?
\end{enumerate}
\section{Simulations}
\label{sec:simulations}

All training and testing data are drawn from the Cosmology and Astrophysics with Machine-learning Simulations \citep[CAMELS;][]{VillaescusaNavarro-2021} project. CAMELS is a suite of hydrodynamical and $N$-body simulations designed to train and test machine-learning models on cosmological data. In the updated $50\,Mpc/h$ suite, each simulation follows \(512^3\) dark-matter particles and, in the hydrodynamical runs, an equal number of gas resolution elements in a periodic comoving box of side length $L = 50\,h^{-1}\,{\rm Mpc}$. The hydrodynamical runs employ the IllustrisTNG galaxy-formation model \citep{weinberger2017supermassive,pillepich2018simulating}, implemented in the \textsc{arepo} moving-mesh code \citep{springel2010arepo}.

\subsection{The SB35 Latin-Hypercube Suite}
\label{subsec:sb35}

We train \textsc{BIND} on the SB35 suite \citep{Genel-2026}, consisting of 1{,}024 paired IllustrisTNG hydrodynamical and $N$-body simulations spanning a 35-dimensional parameter space sampled with a Sobol sequence \citep{Sobol-1967}. The parameter space includes five cosmological parameters and 30 astrophysical subgrid parameters governing stellar feedback, AGN feedback, and related processes in the IllustrisTNG model (for a full breakdown of the parameters see \citealt{Genel-2026}). The Sobol sequence ensures space-filling coverage of this high-dimensional parameter space, with each simulation occupying a unique combination of parameter values and random initial phases.

Each hydrodynamical simulation is paired with a gravity-only $N$-body run evolved from identical initial conditions. Therefore, by holding the initial conditions fixed and barring intrinsic randomness due to butterfly effects \citep{Genel-2019}, differences between the $N$-body and hydrodynamical outputs arise solely from baryonic physics and provide a clean training signal for the mapping $\bm{M}_{\rm DMO} \rightarrow \bm{M}_{\rm hydro}$.

\subsection{Evaluation Suites}
\label{subsec:eval_suites}

We use three independent simulation sets for evaluation.

The \textbf{SB35 test set} consists of the ${\sim}100$ SB35 simulations withheld from training. Each simulation contains a different set of cosmological and astrophysical parameters, as well as different initial conditions.

The \textbf{1P suite} contains simulations in which a single parameter is varied at a time while all others are held at their fiducial values, and all runs share the same initial conditions. 

The \textbf{Cosmic Variance (CV) suite} consists of 27 simulations run at the fiducial IllustrisTNG parameter values but with different random initial conditions.
% TODO(Max, D5): Sec. 5.1-5.2 say "28 CV simulations" and "1,111 CV halos" ---
% reconcile 27 vs 28 here.


%=============================================================
\section{Model and Training}
\label{sec:model}
%=============================================================

Our goal is to learn the conditional distribution
\begin{equation}
    p(\bm{M}_{\rm hydro} \mid \bm{M}_{\rm DMO},\,
    \boldsymbol{\theta}_{\rm TNG},\, z),
    \label{eq:target_distribution}
\end{equation}
where $\bm{M}_{\rm hydro} \in \mathbb{R}^{3 \times H \times W}$ represents the three target projected mass fields of halos in the hydrodynamical simulation (dark matter, gas, and stellar mass), and $\bm{M}_{\rm DMO} \in \mathbb{R}^{1 \times H \times W}$ is the conditioning input, the dark-matter-only patch centered on the halo. The parameter vector $\boldsymbol{\theta}_{\rm TNG} \in \mathbb{R}^{35}$ encodes the full set of cosmological and astrophysical parameters of the simulation (Section~\ref{subsec:sb35}); however, in general, this vector can be any conditioning parameter vector that is associated with the halo (see Appendix~\ref{appendix:obs} for details on a model trained with $\boldsymbol{\theta}\in [Y_{\rm SZ}, T, f_b]$). The scalar $z$ is the redshift of the snapshot from which the halo is drawn: \textsc{BIND} is trained on halos from multiple snapshots and conditions on $z$ alongside $\boldsymbol{\theta}_{\rm TNG}$.
% TODO(Max): specify the training snapshots/redshift range, the z
% normalization, and its embedding in Sec. 3.2/3.3; add the redshift-
% validation subsection promised by the companion paper; thread z through
% the velocity-field signatures of Eqs. (fm_loss)-(euler).

All fields are defined on a regular $H \times W = 128 \times 128$ grid. The central patch spans $6.25\,h^{-1}\,{\rm Mpc}$ at a resolution of $48.8\,h^{-1}\,{\rm kpc\,pixel}^{-1}$.

% [optional --- uncomment if you want to preempt the referee question of why
% no wider environmental conditioning is used:]
% We note that we also trained a variant conditioned on additional
% downsampled maps of the surrounding large-scale structure, and found no
% improvement in any of the validations presented below; the local patch
% and the parameter vector are evidently sufficient for the halo-scale
% fields BIND generates.
The \textsc{BIND} framework is built on a conditional flow matching network that learns the mapping from dark-matter-only halo fields and IllustrisTNG model parameters to the baryonic mass fields of the paired hydrodynamical simulation. In the following subsections we review the flow matching formalism, describe the network architecture, detail the data processing pipeline, and specify the training procedure.

%-------------------------------------------------------------
\subsection{Flow Matching}
\label{subsec:fm}
%-------------------------------------------------------------

The overall goal in our conditional flow matching is to learn a velocity field $\bm{v}_{\phi}$ that transports samples from an isotropic Gaussian prior, $\bm{M}_0 \sim \mathcal{N}(\bm{0},\bm{I})$, to the target baryonic field $\bm{M}_{\rm hydro}$, conditioned on the paired dark-matter-only map $\bm{M}_{\rm DMO}$ and a parameter vector, which in the bulk of this work is the IllustrisTNG model's astrophysical subgrid and the $\Lambda$CDM cosmological parameters $\boldsymbol{\theta}_{\rm TNG}$.

For a uniformly drawn interpolation time $t \sim \mathcal{U}(0,1)$, the intermediate state is the linear interpolation,
\begin{equation}
\begin{split}
    \bm{M}_t &= (1-t)\,\bm{M}_0 + t\,\bm{M}_{\rm hydro}, \\
    \bm{M}_0 &\sim \mathcal{N}(\bm{0},\bm{I}),
\end{split}
    \label{eq:fm_path}
\end{equation}
which defines a straight-line probability path from noise at $t=0$ to data at $t=1$ \citep{lipman2023flowmatching}. The target velocity along this path is the difference between the target and noise fields,
\begin{equation}
    \bm{v}^{\ast} = \bm{M}_{\rm hydro} - \bm{M}_0.
    \label{eq:fm_target}
\end{equation}
We parameterize the velocity field with a neural network $\bm{v}_{\phi}(\bm{M}_t, t, \bm{M}_{\rm DMO}, \boldsymbol{\theta}_{\rm TNG})$, which we write as $\bm{v}_{\phi}$ for brevity. The network is trained to minimize the conditional flow matching objective
\begin{equation}
    \mathcal{L}_{\rm FM}(\phi)
    =
    \mathbb{E}_{t,\,\bm{M}_0,\,\bm{M}_{\rm hydro}}
    \!\left[\,
    \left\|
        \bm{v}_{\phi}
        -
        \left(\bm{M}_{\rm hydro} - \bm{M}_0\right)
    \right\|_2^2
    \right],
    \label{eq:fm_loss}
\end{equation}
where the conditioning inputs $(\bm{M}_{\rm DMO}, \boldsymbol{\theta}_{\rm TNG})$ are held fixed for each training sample.

At inference, we generate baryonic fields by integrating the learned probability-flow ODE from $t=0$ to $t=1$,
\begin{equation}
    \frac{d\bm{M}}{dt}
    =
    \bm{v}_{\phi}\!\left(
        \bm{M},\,t,\,
        \bm{M}_{\rm DMO},\,\boldsymbol{\theta}_{\rm TNG}
    \right),
    \label{eq:ode}
\end{equation}
starting from $\bm{M}(0) \sim \mathcal{N}(\bm{0},\bm{I})$. We integrate Eq.~(\ref{eq:ode}) with a first-order Euler scheme using $N_{\rm steps}=20$ uniform time steps $\Delta t = 1/N_{\rm steps}$,
\begin{equation}
    \bm{M}_{k+1}
    =
    \bm{M}_{k}
    +
    \Delta t\;
    \bm{v}_{\phi}\!\left(
        \bm{M}_k,\,t_k,\,
        \bm{M}_{\rm DMO},\,\boldsymbol{\theta}_{\rm TNG}
    \right).
    \label{eq:euler}
\end{equation}
Because the probability path in Eq.~(\ref{eq:fm_path}) is a straight line, the learned velocity field is nearly constant along each trajectory, and the Euler integrator remains accurate with this modest number of function evaluations \citep{lipman2023flowmatching}.


%-------------------------------------------------------------
\subsection{Data Processing}
\label{subsec:data_processing}
%-------------------------------------------------------------

Converting raw simulation particle data into training samples for \textsc{BIND} involves four main steps: halo identification and centering, random rotation, line-of-sight projection, and field normalization.

We treat the $N$-body L50 CAMELS suite of simulations as the base set from which halos with $M_{200c} \geq 10^{13}\,M_{\odot}\,h^{-1}$\footnote{$M_{200c}$ is the mass enclosed within the radius $R_{200c}$, inside which the mean density is $200\rho_c$, with $\rho_c$ the critical density of the universe.} are identified. For each halo, we exploit the periodic boundary conditions of the simulation box and re-center the volume on the particle with the deepest gravitational potential, placing the halo at $(x,y,z) = (25,25,25)\,{\rm Mpc}\,h^{-1}$.

We place a stationary observer at $(x,y,z) = (25,25,0)\,{\rm Mpc}\,h^{-1}$ and apply a random rotation about the halo center. The periodic boundary conditions ensure that with appropriate tiling of the simulation box, a $50\,{\rm Mpc}\,h^{-1}$ cube remains volume filling after the rotation. The three Euler angles defining each rotation are stored and subsequently applied to the paired hydrodynamical simulation to guarantee exact spatial correspondence between the $N$-body and baryonic projections. We perform this procedure 10 times per halo, drawing independent rotation matrices each time to augment the training set.

Each rotated volume is projected onto a $1024\times1024$ grid along the $z$-axis using cloud-in-cell (CIC) interpolation \citep{hockney}. We retain the central $128\times128$ pixel patch as the primary $N$-body conditioning field $\bm{M}_{\rm DMO}$. For the hydrodynamical simulations, the identical rotation matrices and $N$-body halo centers are used when projecting the dark-matter, gas, and stellar density fields. 

We highlight a few of the design choices described above. First is that this is done in 2-D which allows for architecture exploration and testing at a reduced computational cost. 3D fields increase the memory and computational requirements by orders of magnitude, so this first iteration of BIND remains 2D. % TODO(Max): restore with citation --- source for the claim that 2D
% baryonification performs virtually the same as 3D.
Regardless, extending the procedure to 3D fields is a goal saved for future development of BIND. The other important caveat is the projection depth. Each training sample is $6.25\times6.25\,{\rm Mpc}\,h^{-1}$ and $128^2$ pixels, computed as the projected mass over the full $50\,{\rm Mpc}\,h^{-1}$ line of sight. While the majority of the mass within a halo will be from material associated with that halo, there will of course be foreground and background matter not associated with the halo included in this projection. This choice was made to ease the pasting of the baryonified halos back into $N$-body simulations, which we will discuss in \S~\ref{sec:power_spectrum}. However, we have also trained an equivalent model with projections over 6.25 Mpc/h centered on each halo, and in Appendix~\ref{cube_comp}, we compare the effects of these projections. 

Because projected density fields span several orders of magnitude and contain many near-zero pixels, all map channels are compressed with
\begin{equation}
    \bm{M}_c' = \log_{10}(1 + \bm{M}_c),
    \label{eq:logtransform}
\end{equation}
which maps zero-density pixels to zero while suppressing the extreme dynamic range. Each transformed channel $c$ is then standardized to zero mean and unit variance,
\begin{equation}
    \tilde{M}_c
    = \frac{\bm{M}'_c - \mu_c}{\sigma_c + \epsilon},
    \label{eq:zscore}
\end{equation}
where $\mu_c$ and $\sigma_c$ are estimated from a random subsample of 10{,}000 training maps and $\epsilon = 10^{-8}$ prevents division by zero. This normalization is applied identically to the target baryonic fields, $\bm{M}_{\rm DMO}$, and $\bm{M}_{\rm ls}$.

The parameter vector $\boldsymbol{\theta}_{\rm TNG}$ is normalized component-wise. Parameters spanning many orders of magnitude are first log-transformed before linear rescaling to $[0,1]$:
\begin{equation}
    \tilde{\theta}_j
    =
    \frac{\theta_j^{(\log)} - \theta_{j,\min}^{(\log)}}
         {\theta_{j,\max}^{(\log)} - \theta_{j,\min}^{(\log)} + \epsilon},
    \label{eq:paramnorm}
\end{equation}
where
\begin{equation}
    \theta_j^{(\log)}
    =
    \begin{cases}
        \log_{10}(\theta_j), & \text{if parameter } j \text{ is log-scaled},\\
        \theta_j,            & \text{otherwise.}
    \end{cases}
\end{equation}
The bounds $\theta_{j,\min}$ and $\theta_{j,\max}$ are taken from the SB35 parameter tables \citep{Genel-2026}, so the normalization range is well-defined for any simulation in the suite, not only those in the training subset.

We split the 1{,}024 SB35 simulations into $90\%$ for training and $10\%$ for testing. Within the training set, $80\%$ of halos are used for optimization and $20\%$ for validation. After applying 10 random rotations per halo, the training set contains approximately 360{,}000 projected halo maps spanning the full 35-dimensional parameter space.

%-------------------------------------------------------------
\subsection{Architecture}
\label{subsec:architecture}
%-------------------------------------------------------------

We parameterize the velocity field $\bm{v}_\phi$ with a conditional U-Net \citep{ronneberger2015unet} that receives the concatenated input tensor
\begin{equation}
    \bm{X}^{\rm in}_t
    = \left[\bm{M}_t \;\middle|\; \bm{M}_{\rm DMO}\right]
    \in \mathbb{R}^{(3+1) \times H \times W},
    \label{eq:model_input}
\end{equation}
at each integration time $t$, together with the parameter vector $\boldsymbol{\theta}_{\rm TNG}$. For more details about the architecture, see Appendix~\ref{appendix:architecture}


%-------------------------------------------------------------
\subsection{Training}
\label{subsec:training}
%-------------------------------------------------------------

The model is trained using the AdamW optimizer \citep{loshchilov2019decoupled} with learning rate $\eta_0 = 10^{-4}$ and weight decay $\lambda = 10^{-4}$. Gradients are clipped to a maximum $\ell_2$ norm of 1.0. We train for up to 200 epochs with a batch size of 64 maps per GPU, using PyTorch Lightning's Distributed Data Parallel (DDP) strategy across multiple A100 GPUs in mixed \texttt{bfloat16} precision.

The learning rate follows a two-phase schedule: a linear warm-up from zero to $\eta_0$ over the first $N_{\rm warm}=1{,}000$ gradient steps, followed by cosine annealing to zero over the remainder of training, 
\begin{equation}
    \eta(k) =
    \begin{cases}
        \eta_0\,\dfrac{k}{N_{\rm warm}},
            & k < N_{\rm warm}, \\[8pt]
        \dfrac{\eta_0}{2}
        \!\left[1 + \cos\!\left(
            \pi\,\dfrac{k - N_{\rm warm}}{K_{\rm tot} - N_{\rm warm}}
        \right)\right],
            & k \geq N_{\rm warm},
    \end{cases}
    \label{eq:lrsched}
\end{equation}
where $k$ is the current gradient step and $K_{\rm tot}$ is the total number of steps \citep{loshchilov2017sgdr}.

An exponential moving average (EMA) of the network weights is maintained throughout training with decay $\rho = 0.9999$,
\begin{equation}
    \bar{\phi}_k = \rho\,\bar{\phi}_{k-1} + (1-\rho)\,\phi_k,
    \label{eq:ema}
\end{equation}
and the EMA weights $\bar{\phi}$ are used exclusively at inference \citep{karras2022edm}. Model checkpoints are saved based on validation loss, retaining the three best-performing epochs.

The complete set of training hyperparameters is listed in Table~\ref{tab:training}.

\begin{table}
\centering
\caption{Training hyperparameters.}
\label{tab:training}
\begin{tabular}{lc}
\hline
Hyperparameter & Value \\
\hline
Learning rate $\eta_0$      & $10^{-4}$ \\
Weight decay                & $10^{-4}$ \\
Gradient clip norm          & 1.0 \\
Warm-up steps               & 1{,}000 \\
Max epochs                  & 200 \\
Batch size (per GPU)        & 64 \\
EMA decay $\rho$            & 0.9999 \\
CFG dropout probability     & 0.1 \\
Euler steps at inference    & 20 \\
\hline
\end{tabular}
\end{table}

\section{Fidelity at Fixed Parameters}\label{sec:validation}
We perform a variety of tests to ensure that the trained model generates realistic baryonified halos, defining each validation metric where it is introduced.
% [restructure: the two sentences below are new transition text --- edit voice
% as desired.]
This section addresses the first question posed in the introduction --- whether \textsc{BIND} reproduces the baryonic fields of individual halos at fixed parameters --- progressing from qualitative field-level comparisons to integrated masses, radial profiles, and projected shapes. The response of those fields \emph{across} the parameter space, the second question, is the subject of Section~\ref{sec:param_response}; Table~\ref{tab:validation_summary} indexes the full validation suite of this paper.

% [restructure: new table skeleton --- CHECK numbers and wording.]
\begin{deluxetable*}{llll}
\tabletypesize{\footnotesize}
\tablecaption{Summary of the validation suite.\label{tab:validation_summary}}
\tablewidth{0pt}
\tablehead{\colhead{Test} & \colhead{Suites} & \colhead{Figure(s)} &
           \colhead{Headline result}}
\startdata
Integrated masses & CV, 1P, SB35 & Figs.~\ref{fig:each_halo_mass_comp}--\ref{fig:mass_comp} & unbiased; $<1\%$ DM/total, $\sim1\%$ gas, $\sim5\%$ stars \\
Radial profiles & CV, 1P, SB35 & Fig.~\ref{fig:total_density} & $\lesssim10\%$ at all radii to $3.125\,h^{-1}\,$Mpc \\
Projected shapes & CV, 1P, SB35 & Fig.~\ref{fig:axes_ratios} & $q$ medians and widths matched \\
Mass--parameter correlations & SB35 & Fig.~\ref{fig:mass_corr} & Spearman matched, 15 strongest params \\
Profile--parameter correlations & SB35 & Fig.~\ref{fig:dens_corr} & matched at all radii \\
1P field-level response & 1P & Fig.~\ref{fig:butterfly} & sign and morphology reproduced \\
Redshift dependence & SB35 & Figs.~\ref{fig:z_mass_error}--\ref{fig:z_response} & accuracy flat to $z=2$; $f_b(z)$ evolution matched, conditioning interpolates \\
Baryon fraction & CV, 1P, SB35 & Fig.~\ref{fig:baryon_fraction} & median and scatter, no mass conditioning \\
Scaling relations, residual corner & CV & Figs.~\ref{fig:relations}--\ref{fig:corner} & slopes, normalizations, joint scatter \\
Power spectrum suppression & CV, SB35 & Fig.~\ref{fig:suppression} & tracks the hydro-replace ceiling
\enddata
\end{deluxetable*}

\subsection{Qualitative Field-Level Comparisons}
\label{subsec:qualitative}
%-------------------------------------------------------------

\begin{figure*}
    \centering
    \includegraphics[width=\linewidth]{imgs/flow_matching_diagram.pdf}
    \caption{Flow matching trajectory and field-level comparison for a representative cluster-mass halo ($M_{200c} \sim 10^{14}\,M_{\odot}\,h^{-1}$) drawn from the CV test set, conditioned on the fiducial IllustrisTNG parameter vector $\boldsymbol{\theta}_{\rm fid}$. \textit{Columns 1--3:} Snapshots of the generative trajectory at $t=0$ (Gaussian noise), $t=0.5$ (intermediate state), and $t=1$ (final generated field). \textit{Column 4:} Ground-truth hydrodynamical fields from the paired IllustrisTNG simulation. \textit{Column 5:} Signed residual $(M_{\rm gen} - M_{\rm hydro})/M_{\rm hydro}$. \textit{Rows:} Projected dark-matter, gas, and stellar surface density, respectively. The generated fields reproduce the large-scale morphology of all three components. Residuals in the dark-matter channel are approximately noise-like, while the gas channel retains coherent structure in the residuals, and the stellar channel shows significant pixel-level scatter owing to its sparsity.}
    \label{fig:flow_matching}
\end{figure*}

\begin{figure}
    \centering
    \includegraphics[width=\linewidth]{imgs/fig2_mass_comparison.pdf}
    \caption{Per-halo integrated-mass comparison between the \textsc{BIND}-generated and true hydrodynamical halos for the stellar, gas, and dark matter channels, integrated within each halo's $R_{200c}$ (left column) and over the full $6.25\times6.25\,{\rm Mpc}\,h^{-1}$ patch (right column). Points show individual halos from the SB35 test, CV, and 1P suites; black outlined points show binned medians with $1\sigma$ error bars. Lower panels compare the mass distributions of the true (solid) and generated (outlined) populations.
    % [CHECK caption against the final figure content.]
    }
    \label{fig:each_halo_mass_comp}
\end{figure}

\begin{figure}
    \centering
    \includegraphics[width=\linewidth]{imgs/fig_mass_error.pdf}
    \caption{Relative error in integrated mass, $M_{\rm gen}/M_{\rm hydro}-1$, for each matter component (columns: hydro dark matter, gas, stars, and total) and two apertures (rows: within $R_{200c}$ and within the full $128\times128$ patch). Results are shown independently for the CV (green), 1P (blue), and SB35 test (red) suites. Box edges mark the 25th and 75th percentiles; the central line marks the median. Dashed horizontal lines indicate the zero-bias reference. \textsc{BIND} recovers unbiased integrated masses across all components and test suites, with median errors below $1\%$ for dark matter and total mass, ${\sim}1\%$ for gas, and ${\sim}5\%$ for stars.}
    \label{fig:mass_comp}
\end{figure}


We begin with a qualitative examination of \textsc{BIND}'s outputs for a representative cluster-mass halo ($M_{200c} \sim 10^{14}\,M_{\odot}\, h^{-1}$) drawn from the CV test set. The $N$-body dark-matter field, centered on the halo as described in Section~\ref{subsec:data_processing}, is provided as the conditioning input alongside the fiducial IllustrisTNG parameter vector $\boldsymbol{\theta}_{\rm fid}$.

Fig.~\ref{fig:flow_matching} illustrates both the generative trajectory and the final field-level comparison. Reading left to right, the first three columns show snapshots of the flow at $t=0$ (isotropic Gaussian noise), $t=0.5$ (an intermediate state in which the large-scale morphology is already visible), and $t=1$ (the fully denoised output). The fourth and fifth columns show the ground-truth hydrodynamical fields and the residuals, respectively. By $t=0.5$, the network has already established the large-scale mass distribution and the remaining integration steps primarily sharpen the morphology and populate small-scale features.

Examining the final output, all three baryonic components are reproduced with good visual agreement. The generated dark-matter field closely traces the $N$-body input, and the residuals are approximately noise-like, with no coherent large-scale structure. The gas channel residuals retain more coherent spatial structure, reflecting the sensitivity of the intracluster gas to AGN and stellar feedback processes that are not fully encoded in the dark-matter morphology alone. We see for the gas channels, errors as high as $70\%$, which occur primarily in regions external to the halo cores. The stellar channel, being highly sparse and concentrated in a small number of bright central pixels, shows the largest pixel-level scatter between the generated and true fields. This is expected, as the precise location and amplitude of individual stellar clumps are stochastic at fixed halo mass and parameters, and a generative model is not expected to reproduce a particular realization exactly, but rather the correct statistical ensemble.

These qualitative observations motivate the quantitative analyses that follow. The next sections systematically characterize \textsc{BIND}'s accuracy through integrated mass comparisons (Section~\ref{subsec:integrated_mass}), radially resolved surface density profiles (Section~\ref{subsec:radial_profiles}), and projected shapes (Section~\ref{subsec:shapes}), before turning to the learned dependence on cosmological and astrophysical parameters (Section~\ref{sec:param_response}).

%-------------------------------------------------------------
\subsection{Integrated Mass}
\label{subsec:integrated_mass}
%-------------------------------------------------------------


We begin with the simplest global statistic: the total integrated mass. For every halo with $M_{200c} \geq 10^{13}\,M_{\odot}\,h^{-1}$ in the full test suite, we sum the projected pixel values within $R_{200c}$ of the halo center and within the full $128\times128$ patch for each matter component, and compute the relative error,
\begin{equation}
    \Delta M = \frac{M_{\rm gen}}{M_{\rm hydro}} - 1,
    \label{eq:mass_err}
\end{equation}
where $M_{\rm gen}$ is the generated mass and $M_{\rm hydro}$ is the true mass from the hydrodynamical simulation. 

In Fig.~\ref{fig:each_halo_mass_comp}, we present the match between the stellar, gas and dark matter channels integrated in an aperture of each halo's radius (left column), and with respect to the full $6.25\times6.25\,Mpc/h$ box (right column). All channels for each of the halos from the SB35 test set, CV set and 1P set are shown as points in the upper panels , with binned medians and $1\sigma$ errorbars shown as the black outlined points. The lower panels show each channels mass distribution for the truth (solid) vs the generated. Already, the agreement is spectacular, and highlights that BIND is able to capture the mass distributions of each channel separately. 


Fig.~\ref{fig:mass_comp} shows this relative difference for each matter component, computed independently for each test suite. Across all three suites, the predicted integrated masses are unbiased with small scatter. The median relative error for the hydro dark-matter and total-mass channels is sub-percent, while gas and stellar masses are recovered at the ${\sim}1\%$ and ${\sim}5\%$ level, respectively. Enlarging the aperture from $R_{200c}$ to the full patch (right column) produces little change in the median but noticeably reduces the scatter, consistent with pixel-level noise averaging down over larger areas.

This serves as the most fundamental test of \textsc{BIND}: for virtually all halos across a parameter space entirely unseen during training, the model has learned that total mass is conserved between the $N$-body halo and its hydrodynamical counterpart, and has correctly partitioned that mass among the dark-matter, gas, and stellar components. This implies that the network has internalized the baryon fraction and stellar-to-halo mass relation in some implicit sense, which we investigate further in Section~\ref{sec:results}.


%-------------------------------------------------------------
\subsection{Radial Profiles}
\label{subsec:radial_profiles}
%-------------------------------------------------------------

\begin{figure}[!t]
    \centering
    \includegraphics[width=\linewidth]{imgs/fig_radial_pct_diff.pdf}
    \caption{Mean fractional error in the azimuthally averaged projected surface density profile, $\langle\Sigma_{\rm gen}/ \Sigma_{\rm hydro} - 1\rangle$, as a function of projected radius for each matter component (columns) and test suite (CV: green; 1P: blue; SB35 test: red). Profiles are computed in 15 logarithmically spaced annuli out to $r = 3.125\,h^{-1} \,{\rm Mpc}$. Shaded bands indicate the halo-to-halo scatter (16th--84th percentile). \textsc{BIND} achieves $\lesssim\!10\%$ accuracy at all radii across all three test suites, with the largest deviations confined to the innermost bins where AGN feedback and steep stellar cusps are hardest to reproduce.}
    \label{fig:total_density}
\end{figure}



Rather than integrating within a fixed aperture, we now examine \textsc{BIND}'s accuracy as a function of projected radius. For each halo we compute the azimuthally averaged projected surface density profile in 15 logarithmically spaced annuli extending to the edge of the patch ($r = 3.125\,h^{-1}\,{\rm Mpc}$), and evaluate the mean fractional error
\begin{equation}
    \langle \Delta\Sigma(r) \rangle
    = \left\langle
        \frac{\Sigma_{\rm gen}(r)}{\Sigma_{\rm hydro}(r)} - 1
      \right\rangle,
    \label{eq:profile_err}
\end{equation}
averaged over all halos in each test suite.

Fig.~\ref{fig:total_density} shows this fractional error for each matter component and test suite. \textsc{BIND} achieves $\lesssim\!10\%$ accuracy across all radii and all three evaluation sets. The largest deviations are confined to the innermost radial bins, where the stellar surface density profile is steepest and is the most difficult for the model to learn.



%-------------------------------------------------------------
\subsection{Projected Halo Shapes}
\label{subsec:shapes}
%-------------------------------------------------------------
\begin{figure*}[!t]
    \centering
    \includegraphics[width=\linewidth]{imgs/shape_1_axisratio.pdf}
    \caption{Distributions of the projected axis ratio $q =
             \sqrt{\lambda_{\min}/\lambda_{\max}}$ measured within $R_{200c}$ for halos in the CV (green), 1P (blue), and SB35 test (red) suites. Solid histograms show measurements from the true hydrodynamical fields; outlined histograms show the same quantity measured from \textsc{bind}-generated fields. Columns show the dark-matter, gas, and stellar channels, respectively. The generated distributions match the true medians and widths remarkably well across all channels and test suites. The dark-matter channel peaks near $q\sim 0.8$, the gas channel is more strongly concentrated toward circular shapes, and the stellar channel exhibits the broadest distribution and largest ellipticities.}
    \label{fig:axes_ratios}
\end{figure*}



In this section we explore whether \textsc{bind} also reproduces the \textit{shapes} of the baryonic fields, which are directly relevant to weak-lensing intrinsic alignment studies \citep[e.g.,][]{troxel2015intrinsic, joachimi2015intrinsic}.

For each halo in the test and generated catalogs, we measure the projected 2D shape by computing the intensity-weighted inertia tensor within $R_{200c}$,
\begin{equation}
    Q_{ij} = \frac{\displaystyle\sum_{\rm pix}
                   (r_i - \bar{r}_i)(r_j - \bar{r}_j)\,I}
                  {\displaystyle\sum_{\rm pix} I},
    \label{eq:inertia}
\end{equation}
where $I$ is the projected surface density (mass per pixel) and the sum runs over all pixels within $R_{200c}$, with $R_{200c}$ taken from the $N$-body friends-of-friends catalog. The eigenvalues $\lambda_{\max} \geq \lambda_{\min}$ of $Q_{ij}$ define the projected axis ratio and ellipticity,
\begin{equation}
    q = \sqrt{\frac{\lambda_{\min}}{\lambda_{\max}}}, \qquad
    \varepsilon = \frac{1-q}{1+q},
    \label{eq:shapes}
\end{equation}
where $q=1$ corresponds to a perfectly circular projected profile and $q \to 0$ to a maximally elongated one.

Fig.~\ref{fig:axes_ratios} shows the resulting $q$ distributions for each channel and test suite. The agreement between the true (solid) and generated (outlined) histograms is excellent across all three channels and all three evaluation sets, with the model correctly reproducing both the median and the width of each distribution. A clear ordering emerges among the channels: the dark-matter maps peak near $q \sim 0.8$, reflecting the mildly elliptical projected shapes of cluster-mass halos; the gas channel is more strongly concentrated toward circular configurations, consistent with the tendency of the intracluster medium to thermalize and fill the halo potential isotropically; and the stellar channel shows the broadest distribution with the largest ellipticities, driven by the compact, asymmetric morphologies of brightest cluster galaxies. Notably, the SB35 test set --- which spans the full 35-dimensional parameter space --- shows no appreciable shift in median or spread relative to the CV halos at the fiducial parameters, suggesting that halo shape is determined primarily by the dark-matter structure and is largely insensitive to the subgrid physics parameters over the ranges explored here.

\begin{figure*}[!t]
    \centering
    \includegraphics[width=\linewidth]{imgs/shape_4_spearman.pdf}
    \caption{Spearman rank correlations between the projected shape measures ($q$ and $\varepsilon$) and the astrophysical parameters $\boldsymbol{\theta}_{\rm TNG}$, for the true (filled) and \textsc{BIND}-generated (open) halos, per channel. The correlations are uniformly weak, and \textsc{BIND} reproduces even these weak trends faithfully.
    % [CHECK caption against the final figure content.]
    }
    \label{fig:axes_ratios_correlations}
\end{figure*}

Fig.~\ref{fig:axes_ratios_correlations} confirms this picture quantitatively. The Spearman correlations between $q$ and $\varepsilon$ and the astrophysical parameter vector $\boldsymbol{\theta}_{\rm TNG}$ are uniformly small across all channels, consistent with the visual impression from Fig.~\ref{fig:axes_ratios}. Nevertheless, \textsc{bind} reproduces even these weak correlations faithfully, with close agreement between the true (filled) and generated (open) symbols across all parameters.

We close this subsection with three caveats that should be borne in mind when interpreting these results. First, we measure shapes from projected mass maps, whereas observational shape measurements are made in flux. Because mass and flux are only approximately proportional, a rigorous comparison with observational intrinsic alignment measurements would require modeling the halo spectral energy distribution and computing shapes from mock photometric images. Second, our pixel scale of ${\sim}50\,h^{-1}\,{\rm kpc}$ is comparable to the effective radius of the brightest cluster galaxy at these halo masses, so the central stellar component is effectively unresolved and occupies roughly a single pixel. Third, the projections are performed along the full $50\,h^{-1}\,{\rm Mpc}$ line of sight, so foreground and background mass contributes to both the true and generated images, diluting the intrinsic halo ellipticity signal. Despite these limitations, the excellent agreement between generated and true shape distributions demonstrates that \textsc{bind} has learned the morphological mapping from dark-matter-only halos to their baryonic counterparts --- a non-trivial result that opens the door to large-scale intrinsic alignment emulation.

%-------------------------------------------------------------
%=============================================================
\section{Fidelity across the Parameter Space}
\label{sec:param_response}
%=============================================================
% [restructure: this section intro is new transition text --- edit voice as
% desired. Everything below it is moved verbatim from the former
% Secs. 4.2-4.4.]
The previous section established that \textsc{BIND} reproduces the baryonic fields of individual halos at fixed parameters. We now turn to the second question posed in the introduction: whether the model has learned how those fields \emph{change} across the IllustrisTNG parameter space. We test this at three levels of resolution --- integrated masses, radial profiles, and the fields themselves --- and close with the redshift dependence.

%-------------------------------------------------------------
\subsection{Parameter Response of Integrated Masses}
\label{subsec:mass_response}
%-------------------------------------------------------------

\begin{figure*}
    \centering
    \includegraphics[width=\linewidth]{imgs/fig3_spearman_param_mass.pdf}
    \caption{Spearman rank correlation between integrated mass per component and each of the 15 astrophysical parameters (out of 30) that exhibit the largest correlation with at least one field, evaluated on the SB35 test set. Filled symbols show correlations computed from the true hydrodynamical halos; open symbols show the same quantity computed from \textsc{BIND}-generated fields at identical masses and parameter-space locations. The strong agreement between filled and open symbols demonstrates that the model has correctly internalized the dependence of baryonic mass on the IllustrisTNG subgrid parameters. The stellar channel (stars) exhibits the strongest parameter sensitivity, followed by gas, while the dark-matter channel shows little response to astrophysical parameter variations.}
    \label{fig:mass_corr}
\end{figure*}

A natural follow-up question is whether the model has also learned the \textit{correlation} between integrated mass and individual simulation parameters. We probe this directly in Fig.~\ref{fig:mass_corr} by computing the Spearman rank correlation coefficient between the integrated mass in each component and each element of $\boldsymbol{\theta}_{\rm TNG}$, for both the true and generated halos at identical parameter-space locations. We show the 15 astrophysical parameters (out of 30) that exhibit the largest correlations with any single channel.

The agreement between the true and generated correlations is remarkably close across all fields and parameters. Several physically intuitive trends are also apparent. The stellar mass shows the strongest sensitivity, as it is dominated by a small number of central high-density pixels and is largely decoupled from the gravitational skeleton traced by the $N$-body field. The gas channel shows intermediate sensitivity, as it is correlated with the dark-matter structure but remains responsive to stellar and AGN feedback variations. The dark-matter channel, by contrast, shows little response to the astrophysical parameters, as it is dominated by gravitational forces and directly constrained by the $N$-body input and cosmological parameters.

%-------------------------------------------------------------
\subsection{Parameter Response of Radial Profiles}
\label{subsec:profile_response}
%-------------------------------------------------------------

\begin{figure*}
    \centering
    \includegraphics[width=\linewidth]{imgs/fig3_spearman_profile.pdf}
    \caption{Spearman rank correlation between the azimuthally averaged surface density profile at each radial bin and the 15 astrophysical parameters with the largest correlations, for the true (filled) and \textsc{BIND}-generated (open) halos. The strong agreement between filled and open symbols at all radii demonstrates that the model correctly captures the radially resolved dependence of baryonic structure on subgrid physics. The stellar channel shows the strongest parameter sensitivity in the central regions, and a marked anticorrelation between stellar and gas densities is visible for several feedback parameters.}
    \label{fig:dens_corr}
\end{figure*}

Fig.~\ref{fig:dens_corr} extends the Spearman correlation analysis of Section~\ref{subsec:mass_response} to the radially resolved profiles. For each annular bin and each matter component, we compute the Spearman correlation between the surface density and each element of $\boldsymbol{\theta}_{\rm TNG}$, for both the true and generated halos. The agreement between the two is strong at virtually all radii and in all fields, indicating that \textsc{BIND} has correctly learned the spatially resolved dependence of baryonic structure on the IllustrisTNG subgrid parameters.

As with the integrated mass comparison of Fig.~\ref{fig:mass_corr}, the stellar channel shows the strongest parameter sensitivity, particularly in the central regions of the halo, where the stellar surface density is highest and AGN feedback most directly regulates the star formation rate. There is also a clear anticorrelation between the stellar and gas density responses for several feedback parameters: parameters that drive higher stellar densities tend to suppress the gas density at similar radii, and vice versa. This reflects the competition between gas cooling and star formation on the one hand, and stellar and AGN feedback-driven outflows on the other --- a balance that \textsc{BIND} appears to have internalized from the training data.

%-------------------------------------------------------------
\subsection{Field-Level Parameter Response}
\label{subsec:field_response}
%-------------------------------------------------------------
\begin{figure*}[!t]
    \centering
    \includegraphics[width=\linewidth]{imgs/fig7_1p_field_response.pdf}
    \caption{Field-level parameter response maps for a representative
    cluster-mass halo ($M_{200c} \sim 10^{14}\,M_\odot\,h^{-1}$) from
    the 1P suite. Each block of rows corresponds to a single IllustrisTNG
    subgrid parameter, with the upper row showing the true hydrodynamical
    residual $\delta\bm{M} = \bm{M}_{\rm high} - \bm{M}_{\rm low}$
    (Eq.~\ref{eq:field_residual}) and the lower row showing the
    \textsc{BIND}-generated residual at the same parameter bounds.
    Columns show the dark-matter, gas, and stellar projected surface
    density channels from left to right. The color scale is symmetric
    and identical for true and generated pairs. The dark-matter residuals
    are consistent with noise in both cases. The gas and stellar channels
    reveal physically interpretable structures: \texttt{ASN1} drives
    a positive gas excess in infalling satellites and a stellar suppression
    throughout the halo; AGN parameters (\texttt{AAGN1}, \texttt{AAGN2})
    produce a central gas deficit and peripheral excess consistent with
    AGN-driven outflows. \textsc{BIND} reproduces the sign and spatial
    morphology of these responses in both channels, though with somewhat
    larger pixel-level noise due to the stochasticity of the generative
    sampling.}
    \label{fig:butterfly}
\end{figure*}

One of the hopes for \textsc{BIND} is that unlike radial baryon correction models, it reproduces the mean baryonic fields and also the directional response of those fields to changes in the underlying galaxy-formation parameters. To test this, we exploit the 1P suite, in which a single IllustrisTNG parameter is varied at a time while all others are held at their fiducial values and the initial conditions are kept fixed. For each parameter, we select a halo at the two prior bounds --- the minimum and maximum parameter values --- and compute the signed field-level residual
\begin{equation}\label{eq:field_residual}
    \delta\bm{M}(\bm{x}) = \bm{M}_{\rm high}(\bm{x}) - \bm{M}_{\rm low}(\bm{x}),
\end{equation}
where $\bm{M}_{\rm high}$ and $\bm{M}_{\rm low}$ are the projected mass maps at the upper and lower parameter bounds, respectively. Because the initial conditions are identical across the 1P simulations, any structure in $\delta\bm{M}$ is attributable to baryonic feedback alone rather than cosmic variance. We compute this difference for each of the three output channels (dark matter, gas, and stellar mass) for both the true hydrodynamical fields and the \textsc{BIND}-generated counterparts and show the results in Fig.~\ref{fig:butterfly}.

The left two columns of Fig.~\ref{fig:butterfly} show the dark matter residuals following Eq.~\ref{eq:field_residual} for both the true and the generated. We see that the dark-matter residual maps are mostly consistent with noise with no coherent large-scale structure. This is expected as the dark matter responds mostly gravitationally, and at fixed initial conditions the dark-matter distribution is nearly insensitive to subgrid parameters over the mass and redshift range probed here. 

The middle two panels of Fig.~\ref{fig:butterfly} show the gas residuals, which contains rich, spatially coherent structure varying significantly across the parameters. For example, in ASN1, we see that the cores of both the central halo and exterior halos are condensed as the parameter is turned up, while the exteriors of the halos are suppressed. This is in line with expectations from \citet{Genel-2026} and \citet{Lee-2024}, who have found that an increase in ASN1 starves the black hole, preventing it from ejecting matter deep outside of the central galaxy.
% TODO(Max): confirm the two ASN1-starves-the-black-hole references above
% (Genel-2026 and Lee-2024 are placeholders). NOTE: the commented-out
% paragraphs below this one describe the opposite ASN1 gas-response sign
% (central excess vs. condensed cores) --- reconcile the physics once.
In contrast, when the IMF slope is increased, more massive halos can form, which leads to higher metallicity, cooling, and black hole growth. The core density of the halos are depleted, but the feedback kicks gas deeper outside of the halos shown by the red rings. Similarly, the stellar fields in the right most panels show that changes to the feedback parameters has overall net positive or negative impacts on the total star formation throughout the fields, and that this response is in line with the gas changes. Taken together, this figure demonstrates that \textsc{BIND} has encoded the directional dependence of baryonic processes at the field level.

% Several physically interpretable features are apparent. For \texttt{ASN1}, which governs the wind energy per unit star formation rate in the IllustrisTNG stellar feedback model, the high-minus-low residual gas map shows bright positive regions coinciding with small satellite galaxies and infalling substructure in the outskirts of the halo, as well as a positive excess near the cluster center. Increasing \texttt{ASN1} drives stronger galactic winds that expel gas from the central regions of galaxies, but the compressed and shock-heated gas in infalling satellites is also enhanced, producing the observed pattern of red blotches in the outskirts and a positive central excess when the map at the lower bound is subtracted. The corresponding stellar channel shows a modest but nearly uniform suppression across the map at high \texttt{ASN1}, consistent with stronger stellar feedback quenching star formation throughout the halo population.

% For the AGN feedback parameters (\texttt{AAGN1} and \texttt{AAGN2}), the gas residuals exhibit the opposite sign in the inner regions: the high-minus-low maps show a strong central gas deficit, reflecting the well-known tendency of AGN-driven outflows to evacuate the intracluster medium from halo cores. This is accompanied by a positive gas excess at intermediate projected radii, consistent with the AGN-driven redistribution of gas from the center to the periphery of the halo. These trends are closely mirrored in the \textsc{BIND} outputs. The stellar maps likewise show a mild central suppression at high AGN feedback strength, consistent with AGN quenching of the central galaxy.

\textsc{BIND} recovers the major trends across these parameters with good fidelity. The generated residual maps share the sign, spatial morphology, and approximate amplitude of the true residuals in both the gas and stellar channels, capturing features that span a dynamic range of order $10^{1}$--$10^{2}$ in surface density. The main discrepancy is that the \textsc{BIND} residuals are somewhat noisier than the ground truth: small-scale features in the true maps that arise from the deterministic response of the same halo to a parameter change are partially washed out in the generated fields by the stochastic nature of the sampling. This is not unexpected, as the model generates independent samples from the learned posterior for each parameter value rather than computing a deterministic perturbation; the residual is therefore noisy at the single-sample level even if the mean response is correct. Still, the bulk spatial structure is preserved, and the qualitative parameter responses are clearly recognizable. This provides strong evidence that \textsc{BIND} has internalized the mapping between the IllustrisTNG subgrid parameters and their field-level baryonic signatures, beyond what could be inferred from integrated-mass or azimuthally averaged statistics alone.



%-------------------------------------------------------------
\subsection{Redshift Dependence}
\label{subsec:z_response}
%-------------------------------------------------------------
The results of this section so far have been established at $z = 0$. Because the
principal application of \textsc{BIND} is to paint baryons onto lightcones, which
necessarily span a range of cosmic epochs, we now ask whether the emulator
retains its fidelity --- and, more demandingly, whether it has learned the
\emph{evolution} of the baryonic fields rather than merely their $z=0$
appearance.

\subsubsection{Training across cosmic time}

The model evaluated here is trained on eight snapshots of every SB35 simulation,
spanning $z = 0$ to $z = 4$, listed in Table~\ref{tab:snapshots}. Halos are
identified independently in each snapshot with the same
$M_{200c} \geq 10^{13}\,M_\odot\,h^{-1}$ threshold and the same projection
procedure described in Section~\ref{subsec:data_processing}; the halo catalogs
are therefore not linked across snapshots by a merger tree, and the comparisons
below are between \emph{populations} at each epoch rather than between individual
objects traced through time. Because the halo mass function evolves strongly at
this threshold, the training set is heavily weighted toward low redshift: the
$z = 0$ snapshot contributes $40{,}950$ halos while $z = 4$ contributes only
$89$.

The network conditions on the scale factor $a = 1/(1+z)$ rather than on $z$
itself, since $a$ is bounded on $(0, 1]$ and varies smoothly over the range of
interest. It enters through its own sinusoidal embedding followed by a two-layer
MLP, producing $\bm{e}_a$, which is \emph{summed} with the flow-time and
parameter embeddings before driving the adaptive group normalization in every
residual block (Appendix~\ref{appendix:architecture}, Eq.~\ref{eq:embedding}). The
conditioning is thus continuous in $a$ by construction: nothing in the
architecture restricts the model to the discrete set of training snapshots, a
property we test directly below.

All measurements in this subsection are made on the held-out test simulations,
using $800$ halos per snapshot where available. One subtlety deserves emphasis.
The maps are gridded in \emph{comoving} coordinates, so the projected
$R_{200c}$ aperture must be evaluated at the redshift of the snapshot,
\begin{equation}\label{eq:r200c_z}
    R_{200c}(z) = \left[\frac{3 M_{200c}}{800\pi\,\rho_{c,0}}\right]^{1/3}
                  \frac{1+z}{E(z)^{2/3}} ,
\end{equation}
where $E^2(z) = \Omega_m (1+z)^3 + \Omega_\Lambda$, with $\Omega_m$ taken per
simulation. Using the $z=0$ expression at $z>0$ would
undersize the comoving aperture by $8.5\%$ to $110\%$ over the redshift and
$\Omega_m$ range spanned by SB35, which would by itself manufacture an apparent
evolution in every aperture-integrated quantity. We have verified
Eq.~\ref{eq:r200c_z} against the \texttt{Group\_R\_Crit200} entries of the
friends-of-friends catalogs, which it reproduces to $0.01\%$.

\begin{deluxetable*}{lcccl}
\tabletypesize{\footnotesize}
\tablecaption{Snapshots of the multi-redshift training set. Halo counts are the
number of $M_{200c}\geq10^{13}\,M_\odot\,h^{-1}$ cutouts (one projection per
halo) in each split. The two highest-redshift snapshots are retained in training
but are too sparse to support a percentile estimate, and are excluded from the
evaluation figures.\label{tab:snapshots}}
\tablewidth{0pt}
\tablehead{\colhead{Snapshot} & \colhead{$z$} &
           \colhead{$N_{\rm halo}$ (train)} & \colhead{$N_{\rm halo}$ (test)} &
           \colhead{Role}}
\startdata
90 & 0.000 & 40,950 & 4,272 & train + evaluate \\
82 & 0.209 & 38,371 & 4,009 & train + evaluate \\
74 & 0.468 & 33,482 & 3,490 & train + evaluate \\
60 & 1.045 & 20,383 & 2,099 & train + evaluate \\
52 & 1.484 & 11,904 & 1,245 & train + evaluate \\
44 & 2.002 & 5,608 & 566 & train + evaluate \\
32 & 3.008 & 898 & 95 & train; sparse evaluation \\
24 & 4.008 & 89 & 13 & train only (too sparse to evaluate) \\
\enddata
\end{deluxetable*}

\subsubsection{Accuracy does not degrade with redshift}

Fig.~\ref{fig:z_mass_error} repeats the integrated-mass test of
Fig.~\ref{fig:mass_comp} --- the same estimator, Eq.~\ref{eq:mass_err}, and the
same decomposition into the CV, 1P and SB35 suites --- with redshift on the
abscissa. We report the suites separately rather than pooling them: the suite mix
is not constant across epochs (1P falls from $56\%$ of the halo population at
$z=0$ to $20\%$ at $z=2$, and is absent at $z=0.209$ because those simulations
store no hydrodynamic snapshot there), and since the stellar bias differs in sign
between suites, a pooled median would swing with composition rather than with
model performance.

The dark matter and total mass are accurate and flat at every epoch: within
$R_{200c}$ the dark matter median stays between $+0.3\%$ and $+1.1\%$ and the
total between $+0.2\%$ and $+2.6\%$ in all three suites over
$0 \leq z \leq 2$. The held-out SB35 suite --- the one that tests generalization
across the full parameter space --- is the flattest of the three in every
channel: gas within $+1\%$ and stars between $-8.3\%$ and $-5.7\%$, with the
stellar deficit if anything shrinking toward high redshift. The near-fiducial CV
suite is likewise stable, at $+11\%$ to $+13\%$ in stars and $+4\%$ to $+8\%$
in gas.

The exception is the 1P suite, whose stellar error grows monotonically from
$+9.6\%$ at $z=0$ to $+29\%$ at $z=2$, accompanied by a milder rise in gas. This
is not a mass-selection effect: restricting to a fixed
$10^{13}$--$10^{13.5}\,M_\odot\,h^{-1}$ bin reproduces the same trend
($+10.0\%$ to $+29.1\%$) and the median halo mass of the sample drifts by less
than $0.1$ dex. Two properties of the 1P design plausibly contribute: its
simulations vary a single parameter at a time to the edge of the prior while
holding the rest fiducial, so they populate the axes of the parameter space
rather than its interior; and they share a common set of initial conditions, so
the halos are not independent realizations --- at $z=2$ the $178$ surviving 1P
halos represent far fewer distinct structures than that count suggests. We
therefore regard the high-redshift 1P points as the least constraining of the
three suites, and note that the model's accuracy on the held-out SB35 sample,
which is the relevant one for deployment, does not degrade.

Fig.~\ref{fig:z_profiles} performs the same regrouping on the radially resolved
test of Fig.~\ref{fig:total_density}. The left column retains the comoving
abscissa of that figure, which is the scale on which pixelization and the
over-smoothing of the generative model act; the right column expresses radius in
units of $R_{200c}$, removing the growth of the halo itself from the comparison.
Neither column shows a redshift-dependent inner-radius degradation: the dark
matter profile error stays within $1\%$ and the gas within $2\%$ at all radii and
all epochs, and the curves for the six snapshots overlap within their bootstrap
confidence intervals. The largest coherent effect is a mild ($\sim\!2\%$)
steepening of the innermost gas bin toward high redshift. The stellar profile
error is flat in radius at $-8\%$ to $-14\%$, consistent with the integrated
deficit of Fig.~\ref{fig:z_mass_error} being an amplitude offset rather than a
misplacement of the stellar mass.

\begin{figure*}
    \centering
    \includegraphics[width=\linewidth]{imgs/fig_z_mass_error.pdf}
    \caption{Relative error in integrated mass,
    $M_{\rm gen}/M_{\rm hydro}-1$, as a function of redshift, for each matter
    component and for the total (columns), measured within $R_{200c}$ (top row)
    and over the full $128\times128$ patch (bottom row). This is the estimator
    and the suite decomposition of Fig.~\ref{fig:mass_comp} with redshift on the
    abscissa. Lines and points show the median over the halos of each suite at
    each snapshot; shaded bands show the halo-to-halo 16th--84th percentile
    scatter. The suites are shown separately, not pooled, because their relative
    contribution changes with redshift and their stellar biases differ in sign,
    so a pooled median would track composition rather than performance. The 1P
    curve is absent at $z=0.209$ because those simulations store no
    hydrodynamic snapshot at that epoch. Dark matter and total mass are flat and
    accurate everywhere; the held-out SB35 suite is flat in every channel; the
    growth of the 1P stellar error toward high redshift is discussed in the
    text. Note that the full-patch row is largely insensitive by construction,
    since each generated patch is mass-matched to its truth counterpart.}
    \label{fig:z_mass_error}
\end{figure*}

\begin{figure*}
    \centering
    \includegraphics[width=\linewidth]{imgs/fig_z_radial_pct_diff.pdf}
    \caption{Fractional error in the azimuthally averaged projected surface
    density profile, $\Sigma_{\rm gen}/\Sigma_{\rm hydro}-1$, with curves colored
    by redshift rather than by test suite. \textit{Left:} comoving radius, the
    abscissa of Fig.~\ref{fig:total_density}. \textit{Right:} radius in units of
    $R_{200c}(z)$, which removes the growth of the aperture itself. Rows are the
    three mass channels; note the different vertical scales, which differ by an
    order of magnitude between channels. Shaded bands are bootstrap $68\%$
    confidence intervals on the median, not the halo-to-halo scatter (the latter
    is an order of magnitude larger and is shown in
    Fig.~\ref{fig:z_mass_error}). The redshift curves overlap within their
    uncertainties in every panel: the inner-radius behavior is not
    redshift-dependent.}
    \label{fig:z_profiles}
\end{figure*}

\subsubsection{The model reproduces evolution, not merely stability}

Figs.~\ref{fig:z_mass_error} and~\ref{fig:z_profiles} establish that the error
does not grow with redshift, but they are not by themselves evidence that the
model uses the redshift label at all: an emulator that ignored $z$ entirely and
reproduced each halo from its dark-matter morphology alone could in principle
produce identical figures. The decisive test is a relation whose
\emph{normalization moves with redshift}, which the model can only track if the
conditioning carries physical information.

Fig.~\ref{fig:z_evolution} presents two such relations. The left panel shows the
normalized baryon fraction of Eq.~\ref{eq:f_b} against halo mass at three
redshifts. The relation evolves substantially: in the
$10^{13}$--$10^{13.5}\,M_\odot\,h^{-1}$ bin the true median
$f_b(<\!R_{200c})/f_{b,0}$ rises from $0.872$ at $z=0$ to $0.949$ at $z=2$, a
$+8.9\%$ change, reflecting the shorter time available for feedback to displace
gas beyond the aperture at earlier epochs. \textsc{BIND} recovers $0.846
\rightarrow 0.927$, an evolution of $+9.5\%$, and matches the true median to
better than $2.4\%$ at every redshift and in every mass bin. The middle panel
shows the projected stellar-to-halo mass ratio, which evolves more steeply
($+54\%$ over the same interval in the truth); here the model reproduces the
sense and approximate magnitude of the evolution ($+68\%$) but inherits the
$10$--$17\%$ normalization deficit of the stellar channel. The right panel
isolates the evolution at fixed halo mass, showing that the generated population
tracks the true $f_b(z)$ trajectory in both mass bins.

That the model reproduces the amplitude of the baryon-fraction evolution to
within a percent, while its stellar normalization is biased low by more than ten
percent, is worth stating plainly: the redshift conditioning and the stellar
amplitude are independent failure modes, and only the latter is deficient.

\begin{figure*}
    \centering
    \includegraphics[width=\linewidth]{imgs/fig_z_evolution.pdf}
    \caption{Evolution of two baryonic scaling relations, true versus generated.
    \textit{Left:} normalized baryon fraction $f_b(<\!R_{200c})/f_{b,0}$
    (Eq.~\ref{eq:f_b}) against $M_{200c}$ at three redshifts; solid lines with
    shaded 16th--84th percentile bands are the true hydrodynamical halos, dashed
    lines with open symbols are \textsc{BIND}. \textit{Middle:} the projected
    stellar-to-halo mass ratio for the same halos. Lower sub-panels show the
    ratio of the generated to the true median in each mass bin.
    \textit{Right:} the same baryon fraction as a function of redshift in two
    fixed bins of halo mass, with the shaded band the bootstrap $68\%$
    confidence interval on the true median. Unlike
    Figs.~\ref{fig:z_mass_error}--\ref{fig:z_profiles}, these relations evolve
    strongly with redshift, so reproducing them requires the redshift
    conditioning to carry physical information rather than merely to be
    harmless. Mass bins containing fewer than $40$ halos are not plotted, which
    truncates the massive bin at $z=1.48$.}
    \label{fig:z_evolution}
\end{figure*}

\subsubsection{The redshift conditioning is used, and it interpolates}

Fig.~\ref{fig:z_response} is the redshift analogue of the 1P field-response test
of Fig.~\ref{fig:butterfly}. We hold the dark-matter conditioning patch, the
large-scale context and the full parameter vector $\bm{\theta}_{\rm TNG}$ fixed
at their $z=0$ values, and vary only the conditioning scale factor. The initial
noise of the ODE integration is reseeded identically at every redshift, so a
given halo begins from the same latent state in every panel and any difference in
the output is attributable to the label alone.

The residual maps in the second row show a coherent, spatially structured, and
monotonically growing response: as the label is moved to earlier epochs the
diffuse inter-halo gas is depleted while the halo cores brighten, which is the
expected signature of a less-evolved, more centrally concentrated gas
distribution. The response is not a uniform rescaling --- it has the morphology
of the underlying structure.

The lower row quantifies this over $24$ halos spanning the mass range. Within a
\emph{fixed} aperture, held at $R_{200c}(z=0)$ so that no part of the response is
a trivial consequence of the aperture growing, the enclosed gas mass falls
monotonically to $0.85$ of its $z=0$ value by $z=2$, the stellar mass rises by a
factor $2.2$, and the integrated Compton-$y$ signal rises by a factor $21.6$.
We stress that the fixed aperture here is what makes the change attributable to
the label; it is a different convention from the evolving $R_{200c}(z)$ aperture
used in Fig.~\ref{fig:z_evolution}, and the two need not --- and do not --- move
the baryon fraction in the same direction.

The strongest statement in this figure concerns the points drawn as open
diamonds. These are evaluated at redshifts $z = 0.10, 0.33, 0.70, 1.25, 1.75$,
none of which corresponds to any training snapshot. They fall on the same smooth
monotonic curve as the trained values, with no kinks or steps at the trained
redshifts. The model has therefore learned a continuous function of the scale
factor rather than memorizing a discrete set of epochs --- the property that
licenses its use on lightcones, where every intermediate redshift must be
generated.

\begin{figure*}
    \centering
    \includegraphics[width=\linewidth]{imgs/fig_z_response.pdf}
    \caption{Response of the generated fields to the redshift label at fixed
    dark-matter conditioning. \textit{Top row:} generated gas surface density for
    a single halo ($\log M_{200c} = 14.57$) at four values of the conditioning
    redshift, with the dark-matter patch, the large-scale context, the parameter
    vector and the initial ODE noise all held fixed. \textit{Middle row:} the
    residual $\Delta\log_{10}\Sigma_{\rm gas}$ relative to the $z_{\rm label}=0$
    generation; the leftmost panel is identically zero by construction.
    \textit{Bottom row:} median response over $24$ halos, normalized to the
    $z_{\rm label}=0$ value, measured within a fixed aperture of
    $R_{200c}(z=0)$ so that the change reflects the field response rather than a
    growing aperture; shaded bands are the 16th--84th percentiles over halos.
    Filled circles mark redshifts present in the training set; open diamonds mark
    values ($z = 0.10, 0.33, 0.70, 1.25, 1.75$) that correspond to no training
    snapshot. That the off-grid points lie on the same smooth monotonic curve
    demonstrates that the model has learned a continuous function of the scale
    factor rather than memorizing the discrete training epochs. The
    Compton-$y$ panel is a thermodynamic channel, whose comoving-to-physical
    scale-factor conversions at $z>0$ have not been validated against an
    independent reference; the mass channels carry no such conversion.}
    \label{fig:z_response}
\end{figure*}

\subsubsection{Scope and limitations}

Taken together, these tests establish that \textsc{BIND} generates the baryonic
fields with the same fidelity at $z = 2$ as at $z = 0$, that it reproduces the
evolution of the halo baryon fraction to within a few percent, and that its
redshift conditioning is a smooth, continuously interpolating function of the
scale factor. This validates the emulator across the redshift interval spanned by
the lightcones of the companion paper \citep{Lee-2026c}, whose baryonification
is applied over $z = 0$--$2.5$.

Three limitations should be borne in mind. First, the halo statistics thin
rapidly with redshift at a fixed mass threshold
(Table~\ref{tab:snapshots}); above $z \simeq 2$ the held-out sample is too small
to support the percentile estimates used here, and the $z=4$ snapshot is
excluded from every figure. Extending the validation to higher redshift would
require lowering the mass threshold, which is the same requirement identified in
Section~\ref{subsec:caveats} for the power spectrum. Second, because the halo
catalogs at different snapshots are constructed independently rather than linked
by a merger tree, all statements about evolution here are statements about the
halo population at each epoch, not about the histories of individual objects.
Third, the projections integrate the full $50\,h^{-1}\,{\rm Mpc}$ comoving depth
of the box at every redshift, so aperture-integrated quantities include
line-of-sight contributions from outside the halo; this affects the true and
generated maps identically and so does not bias the comparison, but it does mean
that $f_b(<\!R_{200c})$ as measured here is not directly comparable to a
three-dimensional baryon fraction.

\section{Emergent Joint Structure: the \textsc{BIND} Sandbox}
\label{sec:results}
%-------------------------------------------------------------

In this section we turn to the third question posed in the introduction and explore the extent to which \textsc{BIND} captures correlations between physically distinct baryonic components. We begin by examining the halo baryon fraction across all three test suites to probe whether the model correctly partitions mass among gas, stars, and dark matter as a function of halo mass and cosmic baryon fraction. We then focus on the 1{,}111 CV halos generated at the fiducial IllustrisTNG parameters to investigate a set of inter-component scaling relations: the stellar-to-halo mass relation (SHMR), the baryonic-to-halo mass relation, the gas-to-stellar mass relation, and the dark-matter-to-stellar mass relation. Finally, we decompose the scatter around these relations and show that \textsc{BIND} correctly reproduces the joint distribution of halos that deviate coherently
from multiple relations simultaneously.

%-------------------------------------------------------------
\subsection{Baryon Fractions}
\label{subsec:baryon_fraction}
%-------------------------------------------------------------

For each halo in the three test suites, we compute the projected baryon fraction within $R_{200c}$ as
\begin{equation}\label{eq:f_b}
    f_{b}(<\!r) = \frac{M_{\rm gas}(<\!r) + M_{\rm star}(<\!r)}{M_{\rm total}(<\!r)},
\end{equation}
and normalize by the cosmological baryon fraction $f_{b,0} \equiv \Omega_b/\Omega_m$ appropriate to each simulation. We then bin the halos in logarithmic $M_{200c}$ intervals, computing the median normalized baryon fraction and the halo-to-halo scatter (16th--84th percentile) in each bin, for both the true hydrodynamical halos and the \textsc{BIND}-generated counterparts.

Fig.~\ref{fig:baryon_fraction} shows the resulting $f_b/f_{b,0}$ vs.\ $M_{200c}$ relation for each test suite. In the 1P and SB35 test sets, individual halos are colored by their cosmic baryon fraction $f_{b,0}$, illustrating how the model tracks variations across the full cosmological parameter range. In the CV set --- where all 28 simulations share identical cosmological parameters --- all halos are drawn in a single color. The \textsc{BIND}-generated halos reproduce the median and scatter of the true relation with remarkable fidelity across more than a decade in halo mass and across the diverse parameter coverage of the SB35 suite.

This test is a particularly stringent probe of coherent multi-field generation. The baryon fraction involves the ratio of two separately generated fields --- gas and stellar mass --- to the total projected mass. No halo mass information was supplied as explicit conditioning to the network; nonetheless, the model has emergently learned the halo-mass dependence of the baryon fraction, the sensitivity of $f_b$ to the cosmic baryon fraction $\Omega_b/\Omega_m$, and the correct relative normalization between the gas, stellar, and dark-matter channels simultaneously. The close agreement between the true and generated relations across all three suites therefore reflects a non-trivial internal consistency in the generative model.

\begin{figure}
    \centering
    \includegraphics[width=\linewidth]{imgs/obs_baryon_fraction.png}
    \caption{Normalized baryon fraction $f_b(<\!R_{200c})/f_{b,0}$ as a
    function of halo mass $M_{200c}$ for the three test suites. Points show individual halos; lines and shaded bands show the median and 16th--84th percentile scatter computed in logarithmic mass bins. \textit{True} hydrodynamical halos are shown with filled symbols;  \textsc{BIND}-generated halos with open symbols. In the 1P and SB35  sets, halos are colored by their cosmic baryon fraction $\Omega_b/\Omega_m$; the CV set uses a single color as all simulations share identical cosmological parameters. \textsc{BIND} reproduces both the median trend and scatter of the baryon fraction across the full range of halo masses and parameter values, despite halo mass never being used as a conditioning variable.} \label{fig:baryon_fraction}
\end{figure}

%-------------------------------------------------------------
\subsection{Inter-component Scaling Relations}
\label{subsec:scaling_relations}
%-------------------------------------------------------------

We next investigate whether \textsc{BIND} reproduces the joint covariance structure of the baryonic mass components. For this analysis we focus on the 1{,}111 CV halos drawn from the 28 CV simulations, which provide a clean sample at fixed cosmological and astrophysical parameters. We measure four integrated-mass scaling relations: $M_{\rm gas}$--$M_{\rm star}$, $M_{\rm dm}$--$M_{\rm star}$, the SHMR ($M_{\rm star}$--$M_{200c}$), and the baryonic mass--$M_{200c}$ relation, where all masses are measured within $R_{200c}$. As each relation is well described by a power law over the mass range probed, we fit a simple linear model in log--log space and overlay the best-fit regression on both the true and \textsc{BIND}-generated halo populations.

Fig.~\ref{fig:relations} shows these four panels side by side for the true halos (left) and \textsc{BIND} halos (right). The generated population reproduces not only the slope and normalization of each scaling relation but also the amplitude and asymmetry of the scatter about the best-fit line. Halos that scatter significantly above or below the mean relation in the hydrodynamical simulations have corresponding analogs in the \textsc{BIND} output, indicating that the model has internalized the physical correlations that drive halo-to-halo scatter in baryonic content at fixed mass. This is a non-trivial result: scatter in these relations is driven by the assembly history of each halo and the stochastic interplay of cooling, star formation, and feedback, none of which are directly encoded in the $N$-body input beyond the projected dark-matter morphology.

\begin{figure*}
    \centering
    \includegraphics[width=\linewidth]{imgs/scaling_fits.png}
    \caption{Inter-component mass scaling relations for the 1{,}111 CV halos at the fiducial IllustrisTNG parameters. Each panel shows a different relation in log--log space: gas-to-stellar mass (top left), dark-matter-to-stellar mass (top right), stellar-to-halo mass (SHMR; bottom left), and baryonic mass--$M_{200c}$ (bottom right). In each panel, true hydrodynamical halos are shown as filled circles and \textsc{BIND}-generated halos as open circles. Solid and dashed lines show the best-fit linear regressions to the true and generated populations, respectively, with $\pm 1\sigma$ scatter indicated by shaded bands. \textsc{BIND} accurately reproduces both the slopes and normalizations of all four relations as well as the amplitude and morphology of the scatter, indicating that the model has learned the joint covariance structure of the multi-component baryonic fields.}
    \label{fig:relations}
\end{figure*}

%-------------------------------------------------------------
\subsection{Residual Correlations and the Corner Plot}
\label{subsec:corner}
%-------------------------------------------------------------

To examine the degree to which scatter in different scaling relations is correlated, we identify halos that lie more than $1\sigma$ above or below the best-fit linear relation in each of the four panels of Fig.~\ref{fig:relations}, classifying them as the ``high'' and ``low'' populations for each relation. For each pair of relations we then visualize the joint distribution of the residuals $\Delta\log M_{\rm star}$, $\Delta\log M_{\rm gas}$, $\Delta\log M_{\rm dm}$, and $\Delta\log M_{\rm bar}$ for both the high and low subsamples as a corner plot. We smooth the distributions with a Gaussian KDE and overplot the $1\sigma$, $2\sigma$, and $3\sigma$ contours. The same procedure is applied to the \textsc{BIND}-generated population and overlaid for direct comparison.

Fig.~\ref{fig:corner} presents this corner plot for the true and \textsc{BIND}-generated CV halos. The generated contours reproduce the ground-truth distributions with high fidelity across all pairs of residuals and for both the high and low populations. Several physically interesting features are apparent in this plot. The halos that scatter high in the SHMR also tend to scatter high in the baryonic mass relation, reflecting a correlated excess of both stellar and total baryonic content --- likely associated with particularly efficient early cooling and concentrated mass assembly. The halos that scatter low in the gas-to-stellar mass relation tend to scatter high in the dark-matter-to-stellar mass relation, consistent with AGN-driven gas expulsion redistributing baryons out of the halo while leaving the stellar component largely intact. These correlations span multiple physically distinct processes, making them non-trivial for the network to recover from a projected dark-matter input alone.

The close agreement between the truth and \textsc{BIND} across all of these joint distributions is perhaps the most demanding consistency test presented in this paper. It demonstrates that the model has not merely learned marginal distributions in each channel independently, but has captured the full multivariate covariance structure of the baryonic fields as a function of halo properties. This is a direct consequence of the joint, multi-channel generative architecture of \textsc{BIND}, which produces all three mass fields simultaneously from a single sample of the learned probability path, and therefore naturally preserves inter-channel correlations.

\begin{figure*}[!t]
    \centering
    \includegraphics[width=\linewidth]{imgs/residual_corner_plot.pdf}
    \caption{Corner plot of the residuals $\Delta\log M$ about the best-fit scaling relations of Fig.~\ref{fig:relations} for the CV halo population. Each off-diagonal panel shows the joint distribution of residuals for a pair of relations; diagonal panels show the marginal distributions. Orange (blue) contours correspond to halos that scatter more than $1\sigma$ above (below) the mean scaling relation in any given panel. Filled contours show the true hydrodynamical halos; open contours show the \textsc{BIND}- generated halos, smoothed with a Gaussian KDE. Contour levels correspond to the $1\sigma$, $2\sigma$, and $3\sigma$ enclosed probability. The remarkable agreement between the true and generated contours across all panels demonstrates that \textsc{BIND} has learned the full multivariate covariance structure of the baryonic mass components, including physically motivated correlations such as the coherent excess of stellar and baryonic mass in high-SHMR halos and the anticorrelation between gas and stellar residuals driven by AGN feedback.}
    \label{fig:corner}
\end{figure*}
%-------------------------------------------------------------
\section{Using \textsc{BIND} in the Wild}
\label{sec:power_spectrum}
%-------------------------------------------------------------

One of the primary motivations for \textsc{BIND} --- and the fourth question posed in the introduction --- is to serve as a practical baryonification engine: given an arbitrary $N$-body simulation and a choice of IllustrisTNG galaxy-formation parameters, it should produce a baryonified matter field whose statistical properties match those of a true hydrodynamical run at that parameter point. The key downstream observable we target is the projected matter power spectrum suppression,
\begin{equation}\label{eq:suppression}
    S(k) = \frac{P_{\rm bary}(k)}{P_{\rm DMO}(k)},
\end{equation}
which encodes the net impact of baryonic feedback on the matter distribution as a function of scale. Accurately reproducing $S(k)$ across the full IllustrisTNG parameter space is a prerequisite for using \textsc{BIND} as a baryonic correction tool in cosmological inference pipelines.

%-------------------------------------------------------------
\subsection{BINDing an $N$-body Simulation}
\label{subsec:binding}
%-------------------------------------------------------------

To construct a full baryonified field from an $N$-body simulation, we apply \textsc{BIND} halo by halo and paste the generated patches back into a global projected mass map. The procedure proceeds in four steps.

\begin{enumerate}
    \item \textbf{Halo extraction.} We read in the DMO particle catalog for each CV or SB35 test simulation and identify all halos with $M_{200c} \geq 10^{13}\,M_\odot\,h^{-1}$ using the same halo finder and centering procedure described in Section~\ref{subsec:data_processing}. For each halo, we extract the $128^2$-pixel central patch.

    \item \textbf{BIND generation.} We pass each extracted patch through \textsc{BIND}, conditioned on the cosmological and IllustrisTNG parameter vector $\boldsymbol{\theta}_{\rm TNG}$ of the parent simulation, and generate a single sample of the three-channel baryonic output (dark matter, gas, stellar mass) using 20 Euler integration steps.

    \item \textbf{Replacement and pasting.} The $N$-body DMO projection of the full $50\,h^{-1}\,\mathrm{Mpc}$ box is first projected onto a $1024^2$ grid. Each generated halo patch is then pasted back into the global map at the halo center, replacing the corresponding DMO pixels; each generated field blends into the surrounding DMO field through a raised-cosine taper that begins at $3R_{200c}$ and fades the replacement to zero at the patch edge.
In regions where multiple halo patches overlap, we apply a raised-cosine taper function to blend the edges of adjacent patches, and assign priority to the more massive halo so that the largest halos take precedence in contested regions. This strategy follows the approach of \citet{Lee-2023}, adapted here for \textsc{BIND}-generated rather than analytic profile inputs. The total baryonic map is obtained by summing the gas, stellar, and modified dark-matter channels of the pasted field.
% TODO(Max): confirm the taper onset (3R200c) and functional form match the
% pipeline --- the companion lightcone paper delegates this specification here.
% TODO(Max): confirm the pasting-approach reference (Lee et al. 2023 BCM
% pasting assumed).

    \item \textbf{Power spectrum computation.} We compute the 2D projected matter power spectrum of the full baryonified map using \textsc{Pylians} \citep{Pylians} and compare it to the power spectra of (i) the DMO projection, (ii) the true hydrodynamical map, and (iii) a ``hydro-replace'' control map constructed by pasting the actual hydrodynamical halo patches in place of the \textsc{BIND} outputs using the same cosine-taper procedure. The hydro-replace map isolates the effect of the pasting algorithm itself from any inaccuracy in the \textsc{BIND} generations.
\end{enumerate}

Fig.~\ref{fig:suppression} shows the resulting power spectrum suppression $S(k)$ for each of these maps, along with the ratio of each to the true hydrodynamical power spectrum, for both the CV test suite (top row) and the SB35 test set (bottom row). Several conclusions are immediately apparent.

First, the suppression is incomplete: even the true hydrodynamical power spectrum does not reach the level of suppression reported in full-volume IllustrisTNG simulations at these scales. This is a known consequence of our halo mass threshold of $M_{200c} \geq 10^{13}\,M_\odot\,h^{-1}$. Previous work has shown that recovering the full power spectrum suppression requires resolving halos down to $M_{200c} \sim 10^{12}\,M_\odot$ and accurately capturing the baryonic field out to $\sim 5\,R_{200c}$ \citep{vanDaalen-2020, Lee-2026a}; by restricting to the most massive systems, we are capturing only the dominant contribution to the suppression from cluster-scale feedback. Extending \textsc{BIND} to lower halo masses is a straightforward avenue for future work and would bring the recovered suppression into fuller agreement with the complete hydrodynamical result.

Second, the \textsc{BIND} suppression curve tracks the hydro-replace curve closely across the full range of $k$, indicating that \textsc{BIND} is saturating the performance ceiling set by the pasting algorithm itself rather than being limited by the quality of its generated fields. This is the most important result of this section: for the halos and scales accessible with the current setup, the dominant source of residual error is the halo-pasting procedure and the incompleteness of the halo mass range, not the accuracy of the generative model.

Third, the largest discrepancies between the \textsc{BIND} and true power spectra occur at the smallest scales ($k \gtrsim 5\,h\,\mathrm{Mpc}^{-1}$), where the ratio deviates at the few-percent level with \textsc{BIND} predicting slightly more suppression than the truth. This behavior is consistent with a well-known property of flow matching and diffusion-based generative models: the learned velocity field tends to smooth high-frequency structure because the mean squared error training objective penalizes large-amplitude, small-scale fluctuations more than small-amplitude ones. The result is a mild over-suppression of small-scale power that would be partially mitigated by drawing multiple \textsc{BIND} samples per halo and averaging their contributions to the pasted field. We leave this refinement to future work.

Finally, for the SB35 test suite (which spans the full 35-dimensional parameter space with parameter combinations and initial conditions entirely withheld from training), the median suppression and its halo-to-halo scatter are reproduced faithfully across the ensemble of test simulations, consistent with the parameter-sensitivity results in Section~\ref{subsec:scaling_relations}.

\begin{figure*}
    \centering
    \includegraphics[width=\linewidth]{imgs/fig5b_total_field_pk.pdf}
    \caption{Power spectrum suppression $S(k) = P_{\rm bary}(k)/P_{\rm DMO}(k)$ for the CV test suite (left column) and SB35 test set (right column). \textit{Top panels:} Suppression curves for the true hydrodynamical maps (black), the \textsc{BIND}-baryonified maps (blue), the hydro-replace control (orange dashed), and the DMO reference (gray). Shaded bands indicate the 16th--84th percentile scatter over simulations. \textit{Bottom panels:} Ratio of each baryonified power spectrum to the true hydrodynamical power spectrum, with the DMO-to-true ratio (gray) shown as reference. The \textsc{BIND} curve closely tracks the hydro-replace curve at all scales, indicating that the model saturates the ceiling set by the pasting procedure. Residual deviations at small scales ($k \gtrsim 5\,h\,\mathrm{Mpc}^{-1}$) reflect mild over-smoothing by the generative model and are at the few-percent level. The incomplete suppression relative to the full hydrodynamical result is attributable to the halo mass threshold $M_{200c} \geq 10^{13}\,M_\odot\,h^{-1}$ and the absence of low-mass halos and diffuse gas below this limit.}
    \label{fig:suppression}
\end{figure*}

%-------------------------------------------------------------
\subsection{Implications for the Future}
\label{subsec:implications}
%-------------------------------------------------------------

The results in Section~\ref{subsec:binding} establish that \textsc{BIND} can serve as a practical, fast baryonification engine for large-scale structure analyses. Generating baryonic fields for all halos in a single $50\,h^{-1}\,\mathrm{Mpc}$ simulation takes $\mathcal{O}(\mathrm{minutes})$ on a single GPU --- orders of magnitude faster than re-running the hydrodynamical simulation --- making it feasible to baryonify thousands of $N$-body realizations in the time required to complete a single hydrodynamical run. This speed opens several scientific directions.

The most immediate application is marginalizing over baryonic uncertainty in cosmological inference. By running a suite of $N$-body simulations at different cosmological parameters and baryonifying each with \textsc{BIND} at a range of $\boldsymbol{\theta}_{\rm TNG}$ values, one can construct a joint likelihood over cosmological and galaxy-formation parameters directly from field-level observables such as weak-lensing convergence maps, without relying on analytic profile prescriptions. This is particularly attractive for higher-order statistics --- peak counts, Minkowski functionals, and field-level summaries --- for which analytic baryonic correction tools have been
shown to be inadequate \citep{Lee-2023, Lee-2026a}.

Conversely, \textsc{BIND} can be used in a forward-modeling framework to constrain the IllustrisTNG subgrid parameters from observations. Multi-wavelength data such as X-ray gas fractions, the kinematic Sunyaev-Zel'dovich effect, and optical weak lensing jointly constrain different combinations of the feedback parameters \citep{Kovac-2025}. Because \textsc{BIND} is conditioned on the full 35-dimensional parameter vector, it can in principle propagate observational constraints on these parameters into predictions for any field-level observable, providing a unified framework for multi-probe baryon modeling.

A third direction is the generation of large training sets for summary statistic emulators or simulation-based inference networks. Field-level neural network summaries trained on cosmological maps require thousands of simulation realizations at diverse parameter combinations; \textsc{BIND} can provide these at negligible cost once the $N$-body simulations themselves are available.

A first realization of this program is presented in a companion paper \citep{Lee-2026c}, which deploys \textsc{BIND} --- extended with thermodynamic output channels --- on TNG300-Dark lightcones to build a 314-lightcone feedback atlas of seed-matched convergence, Compton-$y$, and optical-depth maps, measures the response of twelve weak-lensing and gas statistics across the 30-dimensional astrophysical prior, and compresses that response into a small set of measured halo properties.

%-------------------------------------------------------------
\subsection{Caveats with the \textsc{BIND} Approach}
\label{subsec:caveats}
%-------------------------------------------------------------

Despite its encouraging performance, the \textsc{BIND} framework as presented here has several limitations that should be borne in mind when applying it in practice.

\paragraph{Halo mass threshold.} The current model is trained and validated on halos with $M_{200c} \geq 10^{13}\,M_\odot\,h^{-1}$, which contribute the largest individual patches to the baryonified field but represent only a fraction of the halos responsible for the full power spectrum suppression. Extending the training to group- and galaxy-scale halos ($10^{12}$--$10^{13}\,M_\odot$) would require handling a significantly larger number of objects per simulation volume and a wider dynamic range of halo morphologies, but is a natural and necessary next step for achieving percent-level accuracy in $S(k)$ at all relevant scales.

\paragraph{Two-dimensional projections.} \textsc{BIND} operates on projected 2D mass maps rather than 3D particle distributions, which introduces line-of-sight confusion from foreground and background matter unassociated with the target halo. While 2D baryonification can perform comparably to 3D for two-point statistics,% TODO(Max): add the 2D-vs-3D citation (same source as the Sec. 3.2 claim).
the approximation may be less accurate for higher-order statistics that are sensitive to the three-dimensional geometry of halos and filaments. Extending \textsc{BIND} to 3D inputs and outputs is conceptually straightforward but requires substantially more GPU memory and training data.

\paragraph{Training set coverage and simulation suite.} \textsc{BIND} is trained exclusively on the IllustrisTNG galaxy-formation model as implemented in the CAMELS SB35 suite. While this covers a broad, 35-dimensional subgrid parameter space, it is anchored to a single simulation code and a single set of physical prescriptions. Applying \textsc{BIND} to parameter combinations outside the SB35 training range, or to simulation codes with qualitatively different feedback implementations (e.g., SIMBA or EAGLE), would require either retraining on the appropriate simulation suite or demonstrating that the learned mapping generalizes across codes.

\paragraph{Single-sample stochasticity.} At inference, we draw a single \textsc{BIND} sample per halo patch. Because the model is generative, this single sample carries stochastic noise at the pixel level, which is why the power spectra in Fig.~\ref{fig:suppression} show slightly more suppression than the truth at small scales. Drawing multiple samples per halo and averaging before pasting would reduce this stochastic over-smoothing at the cost of a proportional increase in inference time. For applications where small-scale power spectrum accuracy is paramount, this averaging strategy is recommended.

\paragraph{Fixed halo center and orientation.} Each halo patch is projected along a fixed line of sight and pasted back at the halo center identified in the $N$-body catalog. This ignores the orientation dependence of the projected baryonic fields and any mis-centering between the $N$-body potential minimum and the true baryonic center of mass. While these effects are small for cluster-mass halos, they may become relevant for lower-mass systems where baryonic backreaction can shift the halo center appreciably relative to the DMO counterpart.


%-------------------------------------------------------------
\section{Summary and Conclusions}
\label{sec:conclusions}
%-------------------------------------------------------------

We have introduced \textsc{BIND} (Baryonic INpainting with Deep learning), a conditional flow matching model that learns the field-level mapping from dark-matter-only halo environments to their hydrodynamical counterparts, conditioned on the full 35-dimensional cosmological and IllustrisTNG subgrid parameter space of the CAMELS SB35 suite. \textsc{BIND} takes as input a projected $N$-body patch centered on a halo --- together with a parameter vector $\boldsymbol{\theta}_{\rm TNG}$ --- and outputs simultaneous, physically consistent realizations of the dark-matter, gas, and stellar mass fields without imposing any analytic functional form on any component. Applied patch-by-patch to the halos in an $N$-body simulation, \textsc{BIND} produces baryonified cosmological fields at $\mathcal{O}(\mathrm{minute})$ cost per simulation volume on a single GPU, orders of magnitude faster than a hydrodynamical re-simulation.

We validated \textsc{BIND} through a comprehensive suite of tests on three independent evaluation sets --- the SB35 test set, the 1P single-parameter-variation suite, and the CV cosmic-variance suite --- spanning a wide range of halo masses, cosmological parameters, and galaxy-formation physics. Our main findings are the following.

\begin{enumerate}

    \item \textbf{Accurate integrated masses and parameter sensitivity.}
    \textsc{BIND} recovers unbiased integrated masses for all three baryonic components across the full evaluation parameter space, with median relative errors below $1\%$ for dark matter and total mass, $\sim 1\%$ for gas, and $\sim 5\%$ for stars. Spearman rank correlations between the generated integrated masses and the IllustrisTNG subgrid parameters match the ground-truth correlations closely across all channels, demonstrating that the model has correctly internalized the dependence of baryonic content on galaxy-formation physics.

    \item \textbf{Radially resolved profile accuracy.}
    Azimuthally averaged surface density profiles are reproduced to better than $10\%$ accuracy at all projected radii out to $3.125\,h^{-1}\,\mathrm{Mpc}$ across all test suites. The radially resolved Spearman correlations with the IllustrisTNG parameters match the ground truth at virtually all radii, capturing the physically motivated competition between gas cooling and feedback-driven expulsion as a function of projected radius.

    \item \textbf{Projected halo morphology.}
    Distributions of the projected axis ratio $q$ for all three mass channels are reproduced faithfully across all test suites, correctly recovering the ordering of shapes among components: dark matter peaks near $q \approx 0.8$, the intracluster gas is more nearly circular, and stars exhibit the broadest distribution and largest ellipticities. This demonstrates that \textsc{BIND} has learned the anisotropic morphological mapping from DMO to baryonic fields --- a non-trivial capability with direct implications for weak-lensing intrinsic alignment modeling.

    \item \textbf{Field-level parameter response.}
    Using the 1P suite, we showed that \textsc{BIND} correctly reproduces the signed, spatially resolved response of the gas and stellar fields to variations in individual subgrid parameters, including the central gas deficit and peripheral excess driven by AGN feedback (\texttt{AAGN1}, \texttt{AAGN2}) and the enhanced infall gas density and stellar suppression driven by stronger stellar winds (\texttt{ASN1}). The dark-matter channel correctly shows no coherent parameter response, consistent with its gravitational origin.

    \item \textbf{Inter-component correlations and scaling relations.}
    Across the 1{,}111 CV halos at the fiducial parameters, \textsc{BIND} reproduces the slope, normalization, and scatter of all four inter-component scaling relations ($M_{\rm gas}$--$M_{\rm star}$, $M_{\rm dm}$--$M_{\rm star}$, SHMR, baryonic mass--$M_{200c}$). Corner plots of the residuals confirm that the joint multivariate distribution of halos that deviate coherently from multiple relations simultaneously is also recovered, including physically motivated correlations such as the coherent SHMR--baryon excess and the gas--stellar anticorrelation driven by AGN feedback. Halo mass was never used as a conditioning variable, yet the model emergently learned the correct halo-mass dependence of the baryon fraction across more than a decade in $M_{200c}$.

    \item \textbf{Power spectrum baryonification.}
    Baryonifying full $50\,h^{-1}\,\mathrm{Mpc}$ $N$-body projections with \textsc{BIND} reproduces the projected matter power spectrum suppression $S(k) = P_{\rm bary}(k)/P_{\rm DMO}(k)$ to within a few percent across the full range of $k$ probed, for both the CV and SB35 test sets. The \textsc{BIND} suppression curve closely tracks the hydro-replace control, indicating that the model saturates the performance ceiling of the halo-pasting procedure rather than being limited by generative accuracy. Residual over-suppression at small scales ($k \gtrsim 5\,h\,\mathrm{Mpc}^{-1}$) is consistent with the mild high-frequency smoothing inherent to flow matching models trained with an $\ell_2$ objective, and can be mitigated by multi-sample averaging.

\end{enumerate}

Together, these results answer the four questions posed in the introduction. \textsc{BIND} can learn the field-level DMO-to-hydro mapping at the halo level (Q1); it correctly encodes the dependence of baryonic structure on the IllustrisTNG galaxy-formation parameter space (Q2); it captures the inter-component correlations and their parameter dependence (Q3); and it can be deployed as a practical baryonification engine with percent-level accuracy for two-point statistics (Q4).

The present work has several limitations that define a natural roadmap for future development. The most important is the halo mass threshold of $M_{200c} \geq 10^{13}\,M_\odot\,h^{-1}$: extending \textsc{BIND} to group- and galaxy-scale halos ($10^{12}$--$10^{13}\, M_\odot$) would recover the full power spectrum suppression and open the model to a wider range of observational comparisons. \textsc{BIND} currently operates in 2D projection, a pragmatic choice that reduces computational cost and simplifies the halo-pasting step; extending to 3D fields is a natural next iteration. The model is also tied to the IllustrisTNG subgrid framework, and extending it to other galaxy-formation models such as SIMBA or EAGLE would require retraining on the appropriate CAMELS suite, though the architecture and training procedure are directly transferable. Finally, drawing multiple \textsc{BIND} samples per halo and combining them coherently --- either by averaging the pasted fields or by constructing an ensemble power spectrum --- would reduce the stochastic pixel-level noise that currently limits small-scale accuracy.

Looking forward, \textsc{BIND} is well positioned to address some of the most pressing systematics challenges for upcoming weak-lensing surveys. The model's halo-centric, field-level, parameter-conditioned design allows it to serve as a forward-modeling engine for joint constraints on cosmological and galaxy-formation parameters using non-Gaussian statistics --- peak counts, Minkowski functionals, and field-level neural summaries --- for which analytic baryonic correction tools have been shown to be insufficient \citep{Lee-2023,Lee-2026a}. When combined with a fast $N$-body simulator and a flexible posterior sampler, \textsc{BIND} could enable full field-level inference pipelines that marginalize over the IllustrisTNG parameter space at speeds that were previously infeasible, unlocking the full statistical power of surveys such as \textit{Euclid}, the Vera C.\ Rubin Observatory LSST, and the Nancy Grace Roman Space Telescope. The companion paper \citep{Lee-2026c} takes the first of these steps.


%% Please use the acknowledgment and contribution environments. This will 
%% be anonomyized when the "anonymous" style option is used. 
\begin{acknowledgments}
% TODO(Max): funding, computing (Flatiron/CCA clusters, Columbia HPC), the
% CAMELS/Learning-the-Universe teams, and colleague acknowledgments.
\end{acknowledgments}

\begin{contribution}
%%This section gives authors the space to recognize author contributions. The text inside this environment is NOT counted towards the total word quanta. At a minimum, manuscripts are expected to include this text:


%% But authors are expected to provide more specific details, e.g. 
%%
%%SC was responsible for writing and submitting the manuscript.
%%WWM came up with the initial research concept and edited the manuscript.
%%OTS obtained the funding and edited the manuscript.
%%EBF provided the formal analysis and validation. He also edited the manuscript.
%%GEH Supervised the undergraduates, wrote the software and administers the project github and Zenodo repositories.
%%
%% Authors can use the Contributor Role Taxonomy (CRediT) at
%% https://credit.niso.org
%% for ideas on how write a good statement tailored to their needs.

\end{contribution}

%% To help institutions obtain information on the effectiveness of their 
%% telescopes the AAS Journals has created a group of keywords for telescope 
%% facilities.
%
%% Following the acknowledgments section, use the following syntax and the
%% \facility{} or \facilities{} macros to list the keywords of facilities used 
%% in the research for the paper.  Each keyword is check against the master 
%% list during copy editing.  Individual instruments can be provided in 
%% parentheses, after the keyword, but they are not verified.
% \facilities{} --- simulation-based work; no observing facilities.

%% Similar to \facility{}, there is the optional \software command to allow 
%% authors a place to specify which programs were used during the creation of 
%% the manuscript. Authors should list each code and include either a
%% citation or url to the code inside ()s when available.
\software{astropy \citep{2013A&A...558A..33A,2018AJ....156..123A,2022ApJ...935..167A},
          Pylians \citep{Pylians}
          % TODO(Max): add numpy, scipy, matplotlib, PyTorch, PyTorch
          % Lightning, and any halo-finder/analysis codes used.
          }

%% Appendix material should be preceded with a single \appendix command.
%% There should be a \section command for each appendix. Mark appendix
%% subsections with the same markup you use in the main body of the paper.
%%
%% Each Appendix (indicated with \section) will be lettered A, B, C, etc.
%% The equation counter will reset when it encounters the \appendix
%% command and will number appendix equations (A1), (A2), etc. The
%% Figure and Table counter will not reset.

\appendix

\section{Architecture}\label{appendix:architecture}
\subsubsection{Conditioning Embeddings}

The flow time $t$ is encoded with sinusoidal positional embeddings
\citep{dhariwal2021diffusion} of dimension 128, then projected to a
512-dimensional embedding $\bm{e}_t$ via a two-layer MLP with SiLU
activations. The parameter vector $\boldsymbol{\theta}_{\rm TNG}$ is
independently processed by a second two-layer MLP ($35 \to 512 \to
512$, SiLU activations) to yield $\bm{e}_\theta$. For the
redshift-conditioned model the scale factor $a = 1/(1+z)$ is encoded the same
way as the flow time --- a sinusoidal embedding followed by its own two-layer
MLP --- to yield $\bm{e}_a$. The embeddings are summed to form a single
conditioning vector
\begin{equation}
    \bm{e} = \bm{e}_t + \bm{e}_\theta + \bm{e}_a,
    \label{eq:embedding}
\end{equation}
which is broadcast into every residual block of the U-Net. Models trained at a
single epoch simply omit the $\bm{e}_a$ term.

\subsubsection{Encoder--Decoder with Adaptive Group Normalization}

The encoder consists of four resolution stages with channel
multipliers $(1,2,4,8)$ relative to a base width of 128, giving
feature-map widths of 128, 256, 512, and 1024. Each stage contains
two residual blocks, and multi-head self-attention layers (4 heads) are
inserted at resolutions $32^2$ and $16^2$. Spatial downsampling uses
strided convolutions; upsampling uses nearest-neighbor interpolation
followed by convolution. Encoder activations are passed to the decoder
via skip connections by channel-wise concatenation.

Each residual block conditions on $\bm{e}$ through Adaptive Group
Normalization \citep{dhariwal2021diffusion},
\begin{equation}
    {\rm AdaGN}(\bm{h},\bm{e})
    =
    \bigl(1 + \boldsymbol{\gamma}(\bm{e})\bigr)
    \odot {\rm GN}(\bm{h})
    + \boldsymbol{\beta}(\bm{e}),
    \label{eq:adagn}
\end{equation}
where ${\rm GN}$ denotes Group Normalization with 32 groups
\citep{wu2018groupnorm}, and $\boldsymbol{\gamma}$,
$\boldsymbol{\beta}$ are learned linear projections of $\bm{e}$.
This allows the network to modulate its internal representations as a
joint function of flow time and cosmological parameters. The output
projection layer is zero-initialized to stabilize early training.

The full architecture is summarized in Table~\ref{tab:architecture}.

\begin{table}
\centering
\caption{U-Net architecture hyperparameters.}
\label{tab:architecture}
\begin{tabular}{lc}
\hline
Parameter & Value \\
\hline
Base channels              & 128 \\
Channel multipliers        & $(1, 2, 4, 8)$ \\
Residual blocks per stage  & 2 \\
Attention heads            & 4 \\
Attention resolutions      & $32^2,\;16^2$ \\
Embedding dimension        & 512 \\
Dropout                    & 0.1 \\
Output initialization      & Zero \\
\hline
\end{tabular}
\end{table}




\section{Projection-depth comparison}\label{cube_comp}
% [EMPTY] TODO(Max): the comparison of the fiducial 50 Mpc/h-depth
% projections against the 6.25 Mpc/h halo-centered variant promised in
% Sec. 3.2 ("we have also trained an equivalent model...").

\section{An observable-conditioned variant}\label{appendix:obs}
% [EMPTY] TODO(Max): the model conditioned on theta in [Y_SZ, T, f_b]
% promised in Sec. 3. Candidate figures already in imgs/:
% paper_obs_corner.pdf, fb_chain_rule_check.pdf,
% jacobian_correlation_by_mass.pdf. Also unplaced: single_halo_comparison.png;
% fig4_radial_pct_diff.pdf (TODO D4: near-duplicate of fig_radial_pct_diff.pdf
% --- confirm which export is current).

%% For this sample we use BibTeX plus aasjournalv7.bst to generate the
%% the bibliography. The sample7.bib file was populated from ADS. To
%% get the citations to show in the compiled file do the following:
%%
%% pdflatex sample7.tex
%% bibtext sample7
%% pdflatex sample7.tex
%% pdflatex sample7.tex

\bibliography{biblio}{}
\bibliographystyle{aasjournalv7}

%% This command is needed to show the entire author+affiliation list when
%% the collaboration and author truncation commands are used.  It has to
%% go at the end of the manuscript.
%\allauthors

%% Include this line if you are using the \added, \replaced, \deleted
%% commands to see a summary list of all changes at the end of the article.
%\listofchanges

\end{document}

% End of file `sample7.tex'.
